\documentclass[11pt,letterpaper]{article}

% ==============================================================================
%% Encoding & idioma
% ==============================================================================
\usepackage[utf8]{inputenc}
\usepackage[T1]{fontenc}
\usepackage[spanish]{babel}

% ==============================================================================
%% Tipografía y diseño
% ==============================================================================
\usepackage{lmodern}
\usepackage[margin=2.5cm]{geometry}
\usepackage{xcolor}
\usepackage{enumitem}
\usepackage{titlesec}
\usepackage{parskip}
\usepackage{array}
\usepackage{booktabs}
\usepackage{hyperref}

% ==============================================================================
%% Paleta de colores (consistente con plantilla Beamer Colmex)
% ==============================================================================
\definecolor{main}{HTML}{5E002B}
\definecolor{subtitle}{RGB}{75,75,75}
\definecolor{cerulean}{RGB}{0,123,167}
\definecolor{redorange}{RGB}{210,77,49}
\definecolor{forestgreen}{RGB}{36,128,103}

\hypersetup{
    colorlinks=true,
    linkcolor=main,
    citecolor=main,
    urlcolor=cerulean,
    pdftitle={Programación para Proyectos de Datos -- Syllabus},
    pdfauthor={Daniel Kelly}
}

% ==============================================================================
%% Estilo de secciones
% ==============================================================================
\titleformat{\section}
  {\Large\bfseries\color{main}}{\thesection}{1em}{}
\titleformat{\subsection}
  {\large\bfseries\color{main}}{\thesubsection}{1em}{}
\titleformat{\subsubsection}
  {\normalsize\bfseries\color{subtitle}}{}{0em}{}

% ==============================================================================
%% Comandos auxiliares (consistentes con plantilla Beamer)
% ==============================================================================
\newcommand{\blue}[1]{{\color{cerulean}#1}}
\newcommand{\orange}[1]{{\color{redorange}#1}}
\newcommand{\green}[1]{{\color{forestgreen}#1}}

% Marcador de contenido pendiente (para iterar sesión por sesión)
\newcommand{\porDesarrollar}{\textit{\color{redorange}[Por desarrollar]}}

% Bloque estandarizado para cada semana (cada semana contiene 2 sesiones de 1:30 hr)
\newcommand{\semana}[2]{%
  \subsubsection*{\color{main}Semana #1: #2}%
}

% Sub-bloque para cada sesión dentro de una semana
\newcommand{\subsesion}[2]{%
  \medskip\noindent\textbf{\color{main}Sesión #1 --- #2}\par\smallskip%
}

% Referencia a cheatsheets oficiales de Posit (rstudio.github.io/cheatsheets)
% Uso: \cheatsheet{slug-del-archivo}{Título descriptivo}
\newcommand{\cheatsheet}[2]{%
  \href{https://rstudio.github.io/cheatsheets/#1.pdf}{\textit{Cheatsheet: #2}}%
}

% ==============================================================================
%% Bibliografía
% ==============================================================================
\usepackage[style=apa, backend=biber]{biblatex}
\addbibresource{../bibliography.bib}

% ==============================================================================
%% Encabezado
% ==============================================================================
\title{\color{main}\textbf{Programación para Proyectos de Datos}\\[0.3em]
       \large\color{subtitle}Temario del Curso}
\author{Daniel Kelly\\
        \small El Colegio de México --- Licenciatura en Economía}
\date{Otoño 2026}

\begin{document}
\maketitle

% ==============================================================================
\section{Presentación}
% ==============================================================================

El objetivo de este curso es complementar la formación cuantitativa de la Licenciatura en Economía de El Colegio de México, con una visión integral del flujo de un proyecto de datos. Los cursos previos del programa (y los disponibles en otras fuentes) equipan al estudiante con \blue{habilidades específicas} de programación y de análisis estadístico; el énfasis aquí está en cómo esas habilidades se \orange{articulan dentro de un mismo proyecto}: desde estructurar un repositorio reproducible y traer datos de fuentes diversas, hasta limpiarlos, modelarlos y producir un entregable comunicable.

La filosofía del curso es \textbf{\textit{zero-to-hero}}: se asume que el estudiante puede llegar con experiencia previa de programación limitada o nula, y termina capaz de levantar por su cuenta un proyecto de datos completo bajo prácticas reproducibles.

El lenguaje del curso es \textbf{R}. La elección se discute en la primera sesión; en términos generales, R es la lingua franca de la investigación cuantitativa académica y ofrece un ecosistema integrado para todas las etapas de un proyecto de datos. El IDE recomendado para el curso es \textbf{Positron}.

% ==============================================================================
\section{Objetivos}
% ==============================================================================

Al finalizar el curso, el estudiante será capaz de:

\begin{enumerate}[leftmargin=*]
    \item Estructurar un proyecto de datos completo bajo convenciones reproducibles (control de versiones, organización de directorios, gestión de dependencias).
    \item Importar, limpiar y transformar datos provenientes de fuentes heterogéneas (archivos planos, hojas de cálculo, bases de datos relacionales).
    \item Programar funciones reutilizables y aplicar paradigmas de programación funcional para automatizar tareas repetitivas.
    \item Estimar y reportar modelos estadísticos comunes en investigación económica desde el punto de vista programático.
    \item Construir entregables comunicables: visualizaciones, documentos automatizados, aplicaciones interactivas e integraciones con servicios externos.
\end{enumerate}

% ==============================================================================
\section{Prerrequisitos}
% ==============================================================================

\begin{itemize}[leftmargin=*]
    \item No se requiere experiencia previa en programación.
    \item Se asume manejo básico de estadística descriptiva y nociones iniciales de econometría (regresión lineal univariable). Es importante precisar que el curso \textit{no} enseña la teoría de los modelos; solo su implementación programática.
\end{itemize}

% ==============================================================================
\section{Metodología}
% ==============================================================================

\textbf{Duración:} 16 semanas, con 2 sesiones de 1:30 hr por semana (32 sesiones totales).

\textbf{Estructura semanal.} Cada semana tiene dos sesiones con roles diferenciados:

\begin{itemize}[leftmargin=*]
    \item \textbf{\blue{Sesión 1 --- Teoría.}} Exposición conceptual del tema de la semana, con código en vivo a manera de ilustración. 
    \item \textbf{\orange{Sesión 2 --- Práctica.}} Inicia con un \textit{resumen rápido} de los puntos clave de la sesión 1, y se dedica el tiempo restante a \textbf{ejercicios prácticos} resueltos en clase.
\end{itemize}

La proporción típica es de \textbf{$\approx$60\% teoría / 40\% código demostrativo}, aunque puede inclinarse hacia la teoría en semanas más conceptualmente densas.

\textbf{Tareas.} El curso \textit{no} contempla tareas semanales fuera de clase. El trabajo evaluable se concentra en los ejercicios resueltos en la segunda sesión semanal y en el \orange{proyecto integrador} que los estudiantes construyen a lo largo del semestre (ver sección correspondiente).

% ==============================================================================
\section{Evaluación}
% ==============================================================================

La estructura de la evaluación propuesta para el curso atiende a los siguientes porcentajes:

\begin{center}
\begin{tabular}{lr}
\toprule
\textbf{Componente} & \textbf{Peso} \\
\midrule
Checkpoints del proyecto integrador       & 50\% \\
Proyecto final                            & 20\% \\
Ejercicios en clase (segunda sesión semanal)    & 20\% \\
Asistencia                                & 10\%  \\
\bottomrule
\end{tabular}
\end{center}

% ==============================================================================
\section{Estructura del curso}
% ==============================================================================

El curso se organiza en cinco partes a lo largo de 16 semanas:

\smallskip
\noindent\textbf{\textcolor{main}{Parte I --- Fundamentos}}
\begin{itemize}[leftmargin=*, itemsep=1pt]
    \item \textbf{Semana 1.} ¿Por qué R? Setup y primeros pasos.
    \item \textbf{Semana 2.} Estructuras de datos: listas, DataFrames, tibbles, indexación, \texttt{NA} y coerción.
    \item \textbf{Semana 3.} Importación de datos y EDA con base R.
    \item \textbf{Semana 4.} Estructura de proyectos, Git y reproducibilidad.
\end{itemize}

\smallskip
\noindent\textbf{\textcolor{main}{Parte II --- Manipulación}}
\begin{itemize}[leftmargin=*, itemsep=1pt]
    \item \textbf{Semana 5.} \texttt{dplyr} I: \textit{pipe}, verbos básicos y sistema tidyverse.
    \item \textbf{Semana 6.} \texttt{dplyr} II: \textit{joins}, \textit{pivots} y manejo de \texttt{NA}.
    \item \textbf{Semana 7.} Strings, expresiones regulares y \texttt{stringr}.
\end{itemize}

\smallskip
\noindent\textbf{\textcolor{main}{Parte III --- Programación}}
\begin{itemize}[leftmargin=*, itemsep=1pt]
    \item \textbf{Semana 8.} Control de flujo y vectorización.
    \item \textbf{Semana 9.} Funciones: diseño, \textit{scoping} y \textit{debugging}.
    \item \textbf{Semana 10.} Programación funcional: \texttt{purrr} y la familia \texttt{apply}.
\end{itemize}

\smallskip
\noindent\textbf{\textcolor{main}{Parte IV --- Análisis}}
\begin{itemize}[leftmargin=*, itemsep=1pt]
    \item \textbf{Semana 11.} Modelado estadístico programático (\texttt{lm}, \texttt{plm}, \texttt{broom}, \texttt{modelsummary}).
    \item \textbf{Semana 12.} Datos a escala: SQL, \texttt{DBI} y \texttt{dbplyr}.
\end{itemize}

\smallskip
\noindent\textbf{\textcolor{main}{Parte V --- Productos Finales}}
\begin{itemize}[leftmargin=*, itemsep=1pt]
    \item \textbf{Semana 13.} Visualización con \texttt{ggplot2}.
    \item \textbf{Semana 14.} Automatización con \texttt{officer}, \texttt{googledrive} y \texttt{gmailr}.
    \item \textbf{Semana 15.} Aplicaciones interactivas con \texttt{Shiny}.
    \item \textbf{Semana 16.} Interfaces de modelos de lenguaje.
\end{itemize}

% ==============================================================================
\section{Temario detallado}
% ==============================================================================

% ------------------------------------------------------------------------------
\subsection*{\color{main}Parte I --- Fundamentos}
% ------------------------------------------------------------------------------

\semana{1}{¿Por qué R? Setup y primeros pasos}

Esta semana corresponde a la intruducción al curso. Cubre la motivación inicial detrás de elegir R, la instalación y configuración del entorno de trabajo, los primeros conceptos del lenguaje (asignación, vectores, tipos), y las políticas del curso sobre el uso de modelos de lenguaje y/o chatbots.

\textbf{Subtemas:}

\begin{itemize}[leftmargin=*, itemsep=2pt]
    \item \textbf{Bienvenida y motivación del curso.} Estructura del curso y esquema de evaluación.
    \item \textbf{¿Por qué R?} R en el contexto de la investigación económica académica; contraste con Python; razones para descartar Stata. Ecosistema CRAN, tidyverse y Posit.
    \item \textbf{Setup: Positron + R.} Instalación; layout del IDE (consola, source, environment, plots, packages); configuración mínima.
    \item \textbf{Uso responsable de modelos de lenguaje en el curso.} Se expone la diferencia operativa entre usar un LLM como \blue{chatbot} (pegar pregunta, copiar respuesta) y como \orange{agente integrado en el proyecto} (lee archivos, propone cambios, ejecuta comandos con el humano en el \textit{loop} ---por ejemplo, Claude Code de Anthropic o Codex de OpenAI). Depender mal de un LLM \textbf{atrofia la capacidad individual} de razonar sobre código; usado como agente puede acelerar el trabajo de quien ya domina los fundamentos. \textbf{Política del curso:} agentes permitidos en el proyecto integrador; \textit{no} permitidos en los ejercicios de clase. El estudiante debe ser capaz de leer y modificar cualquier línea que entregue.
    \item \textbf{Características del lenguaje.} Interpretado, vectorizado, multiparadigma. Cómo se ejecuta R: REPL, script con \texttt{source()}, ejecutable con \texttt{Rscript}.
    \item \textbf{Primeros objetos en R.} Asignación (\texttt{<-} vs \texttt{=}); tipos atómicos (\texttt{numeric}, \texttt{integer}, \texttt{logical}, \texttt{character}); vectores con \texttt{c()}, vectorización, operaciones aritméticas; funciones de inspección (\texttt{class}, \texttt{typeof}, \texttt{length}, \texttt{str}).
    \item \textbf{Evaluación de condiciones lógicas.} Operadores de comparación (\texttt{==}, \texttt{!=}, \texttt{<}, \texttt{>}, \texttt{<=}, \texttt{>=}) y operadores lógicos (\texttt{\&}, \texttt{|}, \texttt{!}; sus contrapartes escalares \texttt{\&\&}, \texttt{||}). Resultado como vector \texttt{logical}; \textit{recycling} en comparaciones entre vectores. Coerción de \texttt{logical} a numérico (\texttt{TRUE} $\to$ \texttt{1}, \texttt{FALSE} $\to$ \texttt{0}) y su utilidad: \texttt{sum()} sobre un vector lógico cuenta los \texttt{TRUE}, \texttt{mean()} calcula la proporción. Funciones auxiliares (\texttt{any}, \texttt{all}, \texttt{xor}, \texttt{isTRUE}, \texttt{isFALSE}) y manejo de \texttt{NA} en condiciones (la propagación silenciosa es una fuente recurrente de errores).
\end{itemize}

\medskip
\noindent\textbf{Objetivos de aprendizaje}
\begin{enumerate}[leftmargin=*, itemsep=1pt]
    \item Justificar la elección de R sobre alternativas programáticas en el contexto de un proyecto académico de investigación económica.
    \item Tener Positron + R instalados y configurados, capaces de ejecutar un \textit{script} básico desde el IDE y desde la terminal.
    \item Crear, asignar y operar con vectores atómicos en R, identificando su tipo y tamaño.
    \item Distinguir el uso de un LLM como chatbot frente a su uso como agente integrado al proyecto, y la política del curso al respecto.
\end{enumerate}

\medskip
\noindent\textbf{Lecturas}
\begin{itemize}[leftmargin=*, itemsep=1pt]
    \item \textit{R for Data Science} (2.\textsuperscript{a} ed.), \textit{Introduction} (sección preliminar) y Cap.\ 2 \textit{Workflow: basics} \cite{wickham2023r4ds}.
    \item \textit{Advanced R} (2.\textsuperscript{a} ed.), Cap.\ 2 §2.1--2.2 y Cap.\ 3 §3.1--3.2 \cite{wickham2019advr}.
\end{itemize}

\semana{2}{Estructuras de datos e indexación}

Esta semana extiende el vocabulario del lenguaje más allá de los vectores atómicos: introducimos listas, DataFrames y tibbles como las estructuras con las que trabajaremos durante el resto del semestre. Cubrimos también dos mecanismos transversales que aparecen una y otra vez en cualquier proyecto de datos: la \blue{indexación} (cómo se accede a un punto específico de cualquier estructura) y la \blue{coerción de tipos} (qué pasa cuando R convierte automáticamente entre clases). Cerramos con el tratamiento de \orange{valores faltantes}, omnipresentes en datos de encuesta.

\textbf{Subtemas:}

\begin{itemize}[leftmargin=*, itemsep=2pt]
    \item \textbf{Listas.} Contenedores heterogéneos (listas); nombres; acceso con \texttt{\$} y con \texttt{[[]]}; diferencia operativa con vectores atómicos; aplicación típica: agrupar resultados de distintos tipos de un análisis.
    \item \textbf{DataFrames y tibbles.} Estructura por debajo (lista de vectores de igual longitud con clase \texttt{data.frame}); constructores \texttt{data.frame()} y \texttt{tibble()}; inspección (\texttt{names}, \texttt{nrow}, \texttt{ncol}, \texttt{dim}, \texttt{head}, \texttt{str}, \texttt{glimpse}). Diferencias prácticas \texttt{tibble} vs.\ \texttt{data.frame}: impresión, ausencia de \textit{partial matching}, no conversión automática de \textit{strings}, \textit{subsetting} que no simplifica. Recomendación de usar \texttt{tibble} por defecto en el curso.
    \item \textbf{Indexación.} Operadores \texttt{[}, \texttt{[[} y \texttt{\$}, y cuándo usar cada uno (preservar clase vs.\ extraer elemento). Indexación posicional, por nombre, lógica y negativa. La función \texttt{which()}. Asignación por \textit{subsetting}.
    \item \textbf{Coerción de tipos.} Coerción implícita al construir vectores y en operaciones aritméticas; jerarquía \texttt{logical} $\to$ \texttt{integer} $\to$ \texttt{double} $\to$ \texttt{character}. Funciones \texttt{as.numeric()}, \texttt{as.character()}, \texttt{as.logical()} para coerción explícita. Casos comunes en datos de encuesta (variables numéricas guardadas como \textit{string}).
    \item \textbf{Valores faltantes.} Naturaleza del \texttt{NA}; propagación en operaciones; \texttt{is.na()}; argumento \texttt{na.rm = TRUE} en funciones agregadoras. Diferencias entre \texttt{NA}, \texttt{NULL} y \texttt{NaN}, y cuándo aparece cada uno.
\end{itemize}

\medskip
\noindent\textbf{Objetivos de aprendizaje}
\begin{enumerate}[leftmargin=*, itemsep=1pt]
    \item Distinguir entre vectores atómicos, listas, DataFrames y tibbles, y elegir la estructura adecuada para representar un conjunto de datos.
    \item Usar correctamente los operadores \texttt{[}, \texttt{[[} y \texttt{\$} para acceder a elementos según el resultado deseado (preservar clase vs.\ extraer).
    \item Anticipar las consecuencias de la coerción implícita al construir o transformar vectores con elementos de tipos mixtos.
    \item Identificar y manejar valores faltantes (\texttt{NA}) en operaciones, distinguiéndolos de \texttt{NULL} y \texttt{NaN}.
\end{enumerate}

\medskip
\noindent\textbf{Lecturas}
\begin{itemize}[leftmargin=*, itemsep=1pt]
    \item \textit{Advanced R} (2.\textsuperscript{a} ed.), Cap.\ 3 \textit{Vectors} ---§3.2.3 \textit{Missing values}, §3.5 \textit{Lists}, §3.6 \textit{Data frames and tibbles}--- y Cap.\ 4 \textit{Subsetting} (completo) \cite{wickham2019advr}.
    \item \textit{R for Data Science} (2.\textsuperscript{a} ed.), Cap.\ 18 \textit{Missing values} \cite{wickham2023r4ds}.
\end{itemize}

\medskip
\noindent\textbf{Referencia rápida:} \cheatsheet{base-r}{Base R} (Posit) ---resumen de funciones base, indexación y estructuras de datos en una página.

\semana{3}{Importación de datos y EDA con base R}

Esta semana constituye la primera aproximación al trabajo con datos reales. Cubre el flujo inicial de cualquier proyecto: traer los datos a R ---ya sea desde archivos o desde APIs HTTP---, inspeccionar lo que llegó para detectar errores tempranos, calcular estadísticas descriptivas y producir visualizaciones exploratorias. Introduce además el lenguaje del \blue{Análisis Exploratorio de Datos} (EDA) como disciplina.

Hacemos énfasis deliberado en el uso de \orange{funciones base de R} para la exploración: descriptivas y gráficas se hacen con lo que viene incluido en el lenguaje, sin paquetes adicionales. La decisión es pedagógica: el estudiante debe entender qué ofrece R por sí mismo antes de subir al nivel del tidyverse (Semanas 5--7) y de \texttt{ggplot2} (Semana 13).

\textbf{Subtemas:}

\begin{itemize}[leftmargin=*, itemsep=2pt]
    \item \textbf{Importación desde archivos.} \texttt{readr} para CSVs (\texttt{read\_csv}, diferencia con \texttt{read.csv} de base; \texttt{problems()}, \texttt{col\_types}) y \texttt{readxl} para Excel (\texttt{read\_excel}, \texttt{sheets}, \texttt{range}). Argumentos clave: \textit{encoding} (UTF-8 vs.\ Latin-1, frecuente en datos mexicanos), representaciones de \texttt{NA}, \texttt{locale}.
    \item \textbf{Consumo de APIs.} Datos como producto de consumir servicios web. Introducción a \texttt{httr2} como cliente HTTP moderno: construir \textit{requests} con \texttt{request()}, ejecutar con \texttt{req\_perform()}, parsear respuestas JSON con \texttt{resp\_body\_json()}, autenticación por \textit{token} en \textit{headers}. Caso de uso central: la \blue{API SIE de Banco de México} para series macroeconómicas (tipo de cambio, INPC, tasas de interés). Aplicación al proyecto integrador: complementar microdatos del INEGI con contexto macro de BANXICO, habilitando preguntas que cruzan ambas fuentes.
    \item \textbf{Inspección post-importación.} Verificar que la importación funcionó: \texttt{str}, \texttt{glimpse}, \texttt{summary}, \texttt{head}/\texttt{tail}, \texttt{dim}. Errores comunes a detectar: tipos mal inferidos, \texttt{NA}s no reconocidos, \textit{encoding} roto.
    \item \textbf{Estadística descriptiva con base R.} Univariadas: \texttt{mean}, \texttt{median}, \texttt{var}, \texttt{sd}, \texttt{quantile}, \texttt{range}, \texttt{summary}. Frecuencias: \texttt{table}, \texttt{prop.table}. Bivariadas: \texttt{cor}, \texttt{cov}, tablas cruzadas con \texttt{table()}. Manejo de \texttt{NA} en agregaciones (\texttt{na.rm = TRUE}).
    \item \textbf{Visualización exploratoria con base graphics.} \texttt{plot}, \texttt{hist}, \texttt{boxplot}, \texttt{barplot}, \texttt{pairs}. Argumentos básicos (\texttt{main}, \texttt{xlab}, \texttt{ylab}, \texttt{col}, \texttt{pch}). Anotaciones (\texttt{abline}, \texttt{lines}, \texttt{points}, \texttt{text}). Criterios para elegir tipo de gráfica según la pregunta.
    \item \textbf{Flujo de EDA.} Pregunta $\to$ carga $\to$ inspección $\to$ descriptivas $\to$ visualización $\to$ ajuste de pregunta. Documentación de hallazgos. La idea de \orange{iteración rápida} entre código y mirada al dato.
\end{itemize}

\medskip
\noindent\textbf{Objetivos de aprendizaje}
\begin{enumerate}[leftmargin=*, itemsep=1pt]
    \item Importar datos desde archivos (CSV, Excel) verificando \textit{encoding}, tipos y representación de \texttt{NA}.
    \item Consumir APIs HTTP con \texttt{httr2} y parsear respuestas JSON, usando la API SIE de Banxico como caso de estudio.
    \item Realizar una inspección sistemática post-importación que detecte errores tempranos (tipos mal inferidos, \texttt{NA}s no reconocidos).
    \item Calcular estadísticas descriptivas univariadas y bivariadas con funciones de base R, manejando correctamente los \texttt{NA}.
    \item Producir gráficas exploratorias (histograma, \textit{boxplot}, dispersión, barras) con \textit{base graphics}, eligiendo el tipo apropiado según la pregunta.
\end{enumerate}

\medskip
\noindent\textbf{Lecturas}
\begin{itemize}[leftmargin=*, itemsep=1pt]
    \item \textit{R for Data Science} (2.\textsuperscript{a} ed.), Cap.\ 7 \textit{Data import} y Cap.\ 20 \textit{Spreadsheets} \cite{wickham2023r4ds}.
    \item \textit{R for Data Science} (2.\textsuperscript{a} ed.), Cap.\ 10 \textit{Exploratory data analysis} (marco conceptual del EDA; los ejemplos visuales usan \texttt{ggplot2}, que se cubre en la Semana 13) \cite{wickham2023r4ds}.
    \item \textit{An Introduction to R} (manual oficial), Cap.\ 12 \textit{Graphical procedures}, para \textit{base graphics} \cite{rcoreteam-intro}.
\end{itemize}

\medskip
\noindent\textbf{Referencia rápida:} \cheatsheet{data-import}{Data import with \texttt{readr}} (Posit) ---funciones de \texttt{readr} con argumentos clave (\texttt{col\_types}, \texttt{locale}, manejo de \texttt{NA}).

\semana{4}{Estructura de proyectos, Git y reproducibilidad}

Semana que cierra la Parte I del curso. Cubre las herramientas que convierten un \textit{script} suelto en un proyecto reproducible: estructura de directorios, control de versiones con Git y captura de dependencias con \texttt{renv}. La estructura propuesta se basa en la \blue{convención utilizada en proyectos profesionales reales}, simplificada para el contexto académico. La meta práctica al cierre: construir un proyecto versionado en un repositorio, que cualquiera pueda clonar, instalar las dependencias, correr el \textit{script} y obtener exactamente los mismos outputs.

\textbf{Subtemas:}

\begin{itemize}[leftmargin=*, itemsep=2pt]
    \item \textbf{Estructura de un proyecto de datos.} Convención del curso, alineada con la práctica profesional: \texttt{files/} (insumos crudos), \texttt{output/} (datos procesados, gráficas, reportes), \texttt{docs/} (documentación, descriptores de la encuesta) y \texttt{pre/} (scripts de \textit{preprocessing}/ETL). Argumento detrás de la separación: datos, código y documentación viven en lugares distintos por razones distintas (versionado, tamaño, audiencia).
    \item \textbf{Proyectos en Positron (\texttt{.Rproj}).} Función del archivo de proyecto: \textit{working directory} fijado automáticamente al \textit{root} y configuración aislada del proyecto. Mientras los \textit{scripts} se ejecuten desde un proyecto abierto en Positron, los paths relativos funcionan sin herramientas adicionales y las rutas absolutas dejan de tener sentido. Para flujos autónomos (\textit{scripts} corridos desde terminal, \textit{scheduling} externo) existen utilidades específicas; quedan fuera del alcance del curso.
    \item \textbf{Control de versiones con Git.} Comandos básicos: \texttt{init}, \texttt{add}, \texttt{commit}, \texttt{status}, \texttt{log}, \texttt{diff}, \texttt{branch}, \texttt{checkout}. Concepto de \textit{remote}; sincronización con GitHub (\texttt{push}, \texttt{pull}). Convenciones de mensajes de \textit{commit} útiles para uno mismo y para colaboradores.
    \item \textbf{\texttt{.gitignore}: qué NO se versiona.} Datos crudos pesados (\texttt{files/}), credenciales (\texttt{.env}), outputs grandes, archivos del sistema (\texttt{.DS\_Store}), historial de R (\texttt{.Rhistory}, \texttt{.RData}). Por qué versionar datos crudos o credenciales es mala práctica.
    \item \textbf{Reproducibilidad con \texttt{renv}.} \textit{Snapshot} de las dependencias del proyecto: \texttt{renv::init()}, \texttt{renv::snapshot()}, \texttt{renv::restore()}. Cuándo justifica su costo (proyectos colaborativos, entregables académicos) y cuándo es excesivo (\textit{scripts} desechables).
    \item \textbf{README y \textit{master script}.} El \texttt{README.md} como pieza central: descripción del proyecto, cómo correrlo, estructura de directorios, cómo regenerar los outputs. Convención del \textit{master script} (\texttt{pre/master.R}) que orquesta el ETL completo con \texttt{source()} ---patrón heredado de los proyectos profesionales referenciados.
\end{itemize}

\medskip
\noindent\textbf{Objetivos de aprendizaje}
\begin{enumerate}[leftmargin=*, itemsep=1pt]
    \item Diseñar la estructura de directorios de un proyecto de datos siguiendo la convención del curso.
    \item Versionar un proyecto con Git, distinguiendo qué archivos se rastrean y cuáles deben quedar excluidos vía \texttt{.gitignore}.
    \item Usar paths relativos al \textit{root} del proyecto, garantizando que los \textit{scripts} se ejecuten correctamente desde Positron en cualquier máquina.
    \item Capturar y restaurar las dependencias del proyecto con \texttt{renv} para asegurar reproducibilidad.
\end{enumerate}

\medskip
\noindent\textbf{Lecturas}
\begin{itemize}[leftmargin=*, itemsep=1pt]
    \item \textit{R for Data Science} (2.\textsuperscript{a} ed.), Cap.\ 6 \textit{Workflow: scripts and projects} \cite{wickham2023r4ds}.
    \item \textit{Happy Git and GitHub for the useR}, capítulos introductorios sobre \textit{setup} y operaciones básicas con Git \cite{bryan-happygit}.
    \item Documentación oficial de \texttt{renv}, secciones \textit{Introduction} y \textit{Workflow} \cite{ushey-renv}.
\end{itemize}

% ------------------------------------------------------------------------------
\subsection*{\color{main}Parte II --- Manipulación}
% ------------------------------------------------------------------------------

\semana{5}{dplyr I: pipe y verbos básicos}

Esta es la \orange{semana bisagra} del curso. Todo lo que viene a continuación ---limpieza avanzada (S6), programación (S8--S10), modelado (S11) y productos (S13--S16)--- asume fluidez con \texttt{dplyr} y con el \textit{pipe} como mecanismo de composición. La meta no es solo familiaridad sino \blue{fluidez}: al cierre de la semana, el estudiante debe poder traducir cualquier pregunta sobre un tibble en una secuencia de verbos, y leer código tidyverse de terceros sin pausa.

\textbf{Subtemas:}

\begin{itemize}[leftmargin=*, itemsep=2pt]
    \item \textbf{Sistema tidyverse y filosofía.} Qué es el tidyverse y qué paquetes lo componen (\texttt{dplyr}, \texttt{tidyr}, \texttt{ggplot2}, \texttt{readr}, \texttt{tibble}, \texttt{purrr}, \texttt{forcats}, \texttt{stringr}, \texttt{lubridate}). La idea de un \textbf{sistema coherente} ---gramática compartida--- frente a la alternativa de paquetes sueltos. Convención del curso: cargar con \texttt{pacman::p\_load(tidyverse)}.
    \item \textbf{Tidy data.} Las tres reglas (cada variable una columna, cada observación una fila, cada valor una celda) y por qué la mayoría de las herramientas del tidyverse las asumen como precondición. Diagnóstico de formatos no-\textit{tidy} típicos en encuestas (variables-columna con sufijos por año o módulo). El \textit{reshape} hacia \textit{tidy} se introduce conceptualmente; el detalle operativo con \texttt{pivot\_longer()} y \texttt{pivot\_wider()} llega en la Semana 6.
    \item \textbf{El \textit{pipe}.} \textit{Pipe} nativo \texttt{|>} (R 4.1+) y \textit{pipe} del tidyverse \texttt{\%>\%} (\texttt{magrittr}). Semántica: el resultado de la izquierda se inserta como primer argumento de la derecha. Por qué convierte código anidado en código legible. \textbf{Convención del curso:} \textit{pipe} nativo \texttt{|>} por defecto; \texttt{\%>\%} solo cuando se requieren sus capacidades específicas.
    \item \textbf{Los seis verbos de \texttt{dplyr}.} \texttt{filter()} (filas por condición), \texttt{arrange()} (ordenamiento), \texttt{select()} (columnas, con \textit{helpers} \texttt{starts\_with}, \texttt{ends\_with}, \texttt{contains}, \texttt{where}), \texttt{mutate()} (creación/modificación), \texttt{summarize()} (colapso a resumen), y \texttt{group\_by()}/\texttt{ungroup()} (contexto de grupo). \textit{Helpers} frecuentes: \texttt{slice\_*()}, \texttt{count()}, \texttt{distinct()}, \texttt{if\_else()}, \texttt{between()}, \texttt{n()}, \texttt{n\_distinct()}. La idea de los verbos como \textbf{bloques componibles}: cualquier transformación es una secuencia de verbos conectados por el \textit{pipe}.
    \item \textbf{\textit{Non-Standard Evaluation} (mención).} Por qué \texttt{filter(df, x > 5)} funciona sin escribir \texttt{df\$x} ni cotizar \texttt{"x"}. Introducción breve al concepto, suficiente para anticipar las limitaciones que aparecen al escribir funciones que reciben nombres de columnas como argumentos. El tratamiento operativo (\texttt{\{\{~\}\}}, \textit{tunneling}) se difiere a la Semana 9.
\end{itemize}

\medskip
\noindent\textbf{Objetivos de aprendizaje}
\begin{enumerate}[leftmargin=*, itemsep=1pt]
    \item Cargar el ecosistema tidyverse y reconocer qué paquete provee qué funcionalidad.
    \item Identificar si un dataset está en formato \textit{tidy} y, en caso contrario, conceptualizar el \textit{reshape} requerido.
    \item Componer secuencias de verbos \texttt{dplyr} conectados por el \textit{pipe} para resolver preguntas de manipulación de datos.
    \item Leer código tidyverse escrito por terceros sin necesidad de ejecutarlo línea por línea.
\end{enumerate}

\medskip
\noindent\textbf{Lecturas}
\begin{itemize}[leftmargin=*, itemsep=1pt]
    \item \textit{R for Data Science} (2.\textsuperscript{a} ed.), Cap.\ 3 \textit{Data transformation} (completo) \cite{wickham2023r4ds}.
    \item \textit{R for Data Science} (2.\textsuperscript{a} ed.), Cap.\ 5 \textit{Data tidying} ---§5.1 \textit{Introduction} y §5.2 \textit{Tidy data}--- (las secciones §5.3 y §5.4 sobre \textit{pivots} corresponden a la Semana 6) \cite{wickham2023r4ds}.
\end{itemize}

\medskip
\noindent\textbf{Referencia rápida:} \cheatsheet{data-transformation}{Data transformation with \texttt{dplyr}} (Posit) ---los seis verbos y sus \textit{helpers} con ejemplos visuales en una página.

\semana{6}{dplyr II: joins, pivots y manejo de NA}

Esta semana completa el conjunto de herramientas que el estudiante necesita para transformar datos crudos en un dataset \textit{tidy} listo para análisis: cierra el tema de \textit{tidy data} iniciado en la Semana 5 con los \textit{pivots}, introduce los \textit{joins} entre tablas, formaliza el manejo avanzado de \texttt{NA}, y cubre la construcción de variables categóricas derivadas. Al cierre de la semana, el estudiante tiene casi todo lo necesario para armar el pipeline E$\to$T del flujo planteado al inicio del curso (la Semana 7 añade \textit{strings} y \textit{regex}, también frecuentes en la limpieza).

\textbf{Subtemas:}

\begin{itemize}[leftmargin=*, itemsep=2pt]
    \item \textbf{Pivots.} \texttt{pivot\_longer()} y \texttt{pivot\_wider()} para \textit{reshape} entre formato ancho y largo. Argumentos clave: \texttt{cols}, \texttt{names\_to}, \texttt{values\_to} (\textit{longer}); \texttt{names\_from}, \texttt{values\_from} (\textit{wider}). Casos típicos en encuestas: variables/columnas con sufijos por año o módulo (ENSU, ENIGH). 
    \item \textbf{Joins entre tablas.} \textit{Mutating joins} (\texttt{left\_join}, \texttt{right\_join}, \texttt{inner\_join}, \texttt{full\_join}) frente a \textit{filtering joins} (\texttt{semi\_join}, \texttt{anti\_join}). Concepto de \textit{key} y sintaxis preferida con \texttt{join\_by()}. Errores comunes: cardinalidad inesperada (\textit{many-to-many}), \textit{keys} duplicadas, \textit{joins} silenciosos sobre columnas mal pareadas. Como mitigación introducimos \blue{\texttt{tidylog}}, que intercepta operaciones de \texttt{dplyr} y \texttt{tidyr} y reporta en consola cuántas filas se afectaron, cuántas no encontraron \textit{match} y cuándo aparecen \texttt{NA} nuevos ---es decir, convierte fallas silenciosas en mensajes explícitos. Aunque se introduce con los \textit{joins}, aplica a todos los verbos vistos esta semana y la anterior.
    \item \textbf{Manejo avanzado de \texttt{NA}.} \texttt{replace\_na()} para rellenar valores específicos; \texttt{coalesce()} para tomar el primer no-\texttt{NA} entre varias columnas; \texttt{drop\_na()} para descartar filas con \texttt{NA} en columnas indicadas. \blue{Decisión deliberada}: cuándo conviene propagar el \texttt{NA} y cuándo intervenir, y por qué la elección debe ser explícita y documentada.
    \item \textbf{Construcción de variables categóricas.} \texttt{if\_else()} vectorizado para condiciones binarias y \texttt{case\_when()} para reglas múltiples mutuamente excluyentes. Mención a \texttt{forcats} para reordenar y recodificar factores (sin profundizar; aparece según necesidad en el proyecto en el que se trabaje).
\end{itemize}

\medskip
\noindent\textbf{Objetivos de aprendizaje}
\begin{enumerate}[leftmargin=*, itemsep=1pt]
    \item Hacer \textit{reshape} de un dataset entre formatos ancho y largo según la operación que se quiera realizar.
    \item Combinar múltiples tablas con los \textit{joins} apropiados, identificando \textit{keys} y anticipando problemas de cardinalidad.
    \item Tomar decisiones explícitas sobre el manejo de \texttt{NA} durante la limpieza, distinguiendo cuándo propagarlos y cuándo intervenir.
    \item Construir variables derivadas categóricas a partir de reglas explícitas con \texttt{if\_else()} y \texttt{case\_when()}.
\end{enumerate}

\medskip
\noindent\textbf{Lecturas}
\begin{itemize}[leftmargin=*, itemsep=1pt]
    \item \textit{R for Data Science} (2.\textsuperscript{a} ed.), Cap.\ 5 \textit{Data tidying} ---§5.3 \textit{Lengthening data} y §5.4 \textit{Widening data} \cite{wickham2023r4ds}.
    \item \textit{R for Data Science} (2.\textsuperscript{a} ed.), Cap.\ 19 \textit{Joins} (completo) \cite{wickham2023r4ds}.
    \item \textit{R for Data Science} (2.\textsuperscript{a} ed.), Cap.\ 18 \textit{Missing values} ---revisión operativa de \texttt{replace\_na()}, \texttt{coalesce()} y \texttt{drop\_na()} \cite{wickham2023r4ds}.
\end{itemize}

\medskip
\noindent\textbf{Referencia rápida:} \cheatsheet{tidyr}{Data tidying with \texttt{tidyr}} (Posit) ---\textit{pivots}, manejo de \texttt{NA} y operaciones de \textit{tidying} en una página.

\semana{7}{Strings, expresiones regulares y \texttt{stringr}}

Esta semana cierra la Parte II del curso. La manipulación de texto es omnipresente en la limpieza de datos: catálogos con códigos pegados a etiquetas, espacios extra, capitalización inconsistente, \textit{paths} de archivos que cambian por trimestre o año. Cubrimos \texttt{stringr} como sistema unificado para operaciones sobre \textit{strings}, las expresiones regulares como lenguaje para describir patrones, y \texttt{glue} para construir \textit{strings} dinámicamente. 

\textbf{Subtemas:}

\begin{itemize}[leftmargin=*, itemsep=2pt]
    \item \textbf{Strings en \textit{base R}: lo mínimo.} \texttt{paste()}/\texttt{paste0()}, \texttt{nchar()}, \texttt{tolower()}/\texttt{toupper()}, \texttt{substr()}. Suficiente para tareas triviales; insuficiente para limpieza sistemática de microdatos.
    \item \textbf{\texttt{stringr}: Flujo y vocabulario consistente.} El patrón \texttt{str\_*()} ofrece un sistema coherente análogo al de \texttt{dplyr}. Funciones \textit{core}: detección (\texttt{str\_detect}), extracción (\texttt{str\_extract}, \texttt{str\_match}), reemplazo (\texttt{str\_replace}, \texttt{str\_replace\_all}), división (\texttt{str\_split}), \textit{recortado} (\texttt{str\_trim}, \texttt{str\_squish}), \textit{padding} (\texttt{str\_pad}), capitalización (\texttt{str\_to\_lower}, \texttt{str\_to\_upper}, \texttt{str\_to\_title}), longitud (\texttt{str\_length}).
    \item \textbf{Expresiones regulares.} Caracteres literales y metacaracteres. Clases de caracteres (\texttt{[abc]}, \texttt{[a-z]}, \texttt{\textbackslash d}, \texttt{\textbackslash s}, \texttt{\textbackslash w}). Cuantificadores (\texttt{*}, \texttt{+}, \texttt{?}, \texttt{\{n,m\}}). \textit{Anchors} (\texttt{\textasciicircum}, \texttt{\$}). Grupos de captura (\texttt{()}) y referencias hacia atrás (\texttt{\textbackslash 1}). Casos típicos: validación de formato (CURP, RFC) y extracción de \textit{tokens}.
    \item \textbf{\texttt{glue} para construcción dinámica de \textit{strings}.} Interpolación de variables (\texttt{glue("Hola \{nombre\}")}), iteración natural sobre vectores, formateo \textit{inline}. Comparación con \texttt{paste0()} y \texttt{sprintf()}; cuándo conviene cada uno. Aplicación a generación de \textit{paths}, mensajes de \textit{log} y plantillas de reporte.
    \item \textbf{Documentar el flujo con \texttt{log()}.} Más allá del \texttt{print()} ocasional, conviene incorporar al pipeline mensajes estructurados que cuenten qué está pasando: cuántas filas se leyeron, cuántas se descartaron, cuánto tardó cada paso. \blue{Funciones de \textit{base R}} para emitir mensajes: \texttt{message()} (canal \textit{stderr}, suprimible), \texttt{warning()} y \texttt{stop()} (señalización jerarquizada de problemas). El paquete \texttt{logger} ofrece un sistema con niveles (\texttt{log\_info}, \texttt{log\_warn}, \texttt{log\_error}, \texttt{log\_debug}), \textit{timestamps} automáticos y \textit{appenders} configurables (consola, archivo, ambos). \blue{Patrón típico}: cada función del pipeline emite un \texttt{log\_info()} con \texttt{glue} al inicio y al cierre, dejando un rastro auditable de la corrida en \texttt{output/log.txt}; cuando algo falla días después, el \textit{log} dice exactamente dónde se rompió y con qué insumos.
    \item \textbf{Texto como dato: distancias y vectorización.} Más allá de tratar \textit{strings} como secuencias de caracteres, existe un terreno donde el texto se manipula como objeto matemático. Introducción al paquete \texttt{stringdist} para distancias entre cadenas (Levenshtein, Jaro-Winkler, Hamming) ---herramienta fundamental para \textit{fuzzy matching} y deduplicación, p.\ ej.\ al cruzar nombres de municipios entre dos fuentes con variaciones ortográficas. Mención somera a \texttt{tidytext} para tokenización y análisis de frecuencias, y a la idea de \textit{word embeddings} como representación vectorial densa (sin profundizar: el tratamiento serio pertenece a un curso de NLP).
\end{itemize}

\medskip
\noindent\textbf{Objetivos de aprendizaje}
\begin{enumerate}[leftmargin=*, itemsep=1pt]
    \item Manipular vectores de \textit{strings} con \texttt{stringr}, distinguiendo cuándo basta una operación de \textit{base R} y cuándo se requiere el sistema.
    \item Leer y escribir expresiones regulares para detectar, extraer y reemplazar patrones en datos textuales.
    \item Construir \textit{strings} dinámicos con \texttt{glue}, distinguiéndolos de las alternativas de \textit{base R}.
    \item Limpiar inconsistencias textuales típicas en datos crudos (capitalización, espacios extra, formatos heterogéneos) combinando \texttt{stringr} con \textit{regex}.
    \item Reconocer cuándo el texto debe tratarse como secuencia de caracteres frente a cuándo conviene tratarlo como objeto matemático (\texttt{stringdist} y técnicas relacionadas), y aplicar \textit{fuzzy matching} a problemas de deduplicación.
\end{enumerate}

\medskip
\noindent\textbf{Lecturas}
\begin{itemize}[leftmargin=*, itemsep=1pt]
    \item \textit{R for Data Science} (2.\textsuperscript{a} ed.), Cap.\ 14 \textit{Strings} (completo) \cite{wickham2023r4ds}.
    \item \textit{R for Data Science} (2.\textsuperscript{a} ed.), Cap.\ 15 \textit{Regular expressions} (completo) \cite{wickham2023r4ds}.
    \item \textbf{Opcional:} \textit{Text Mining with R: A Tidy Approach} (Silge \& Robinson) ---para quienes quieran profundizar en el tratamiento de texto como dato \cite{silge-tidytext}.
\end{itemize}

\medskip
\noindent\textbf{Referencia rápida:} \cheatsheet{strings}{String manipulation with \texttt{stringr}} (Posit) ---API de \texttt{stringr} con ejemplos visuales; \cheatsheet{regex}{Regular Expressions} (Posit) ---tabla compacta de metacaracteres, cuantificadores y \textit{anchors}.

% ------------------------------------------------------------------------------
\subsection*{\color{main}Parte III --- Programación}
% ------------------------------------------------------------------------------

\semana{8}{Control de flujo y vectorización}

Esta semana abre la Parte III del curso (Programación). Cubre la sintaxis de control de flujo en R ---condicionales y \textit{loops}--- y, sobre todo, la \blue{vectorización} como alternativa que el lenguaje ofrece por defecto. La meta pedagógica es que el estudiante reconozca cuándo un \textit{loop} es genuinamente necesario y cuándo es un parche por no haber pensado en términos de vectores. Es preparación directa para las dos semanas siguientes: funciones (S9) y programación funcional con \texttt{purrr} (S10).

\textbf{Subtemas:}

\begin{itemize}[leftmargin=*, itemsep=2pt]
    \item \textbf{Condicionales: \texttt{if}, \texttt{else}, \texttt{switch}.} Sintaxis básica. \texttt{if} opera sobre escalares; con un vector lanza \textit{warning} o solo evalúa el primer elemento. Uso de \texttt{switch()} para múltiples ramas mutuamente excluyentes. Distinción importante: \texttt{if} es para flujo de programa; \texttt{if\_else()} y \texttt{case\_when()} (vistos en S6) son para construir variables vectorizadas dentro de \texttt{mutate()}.
    \item \textbf{\textit{Loops}: \texttt{for} y \texttt{while}.} Sintaxis e iteración. Iteración segura sobre secuencias con \texttt{seq\_along()} y \texttt{seq\_len()}. Iteración sobre listas. \blue{Pre-asignación} de la salida como buena práctica; el anti-patrón de crecer un vector dentro del \textit{loop}. Control fino con \texttt{break} y \texttt{next}. Casos donde el \textit{loop} es genuinamente necesario: operaciones dependientes del estado previo, \textit{side effects} (lectura/escritura de archivos), ETL secuencial sobre años o trimestres.
    \item \textbf{Vectorización como \textit{default}.} R es un lenguaje vectorizado: la mayoría de operaciones aritméticas, lógicas y funciones aplican elemento a elemento sobre vectores. Comparación lado a lado de código con \textit{loop} frente a código vectorizado, en claridad y en rendimiento. Reglas de \textit{recycling}. Funciones acumuladoras (\texttt{cumsum}, \texttt{cummax}, \texttt{cumprod}, \texttt{cumany}, \texttt{cumall}).
    \item \textbf{Patrones y anti-patrones.} \textit{Vector growing} (\texttt{x <- c(x, ...)} dentro de \textit{loop}) frente a pre-asignación. Iteración por valores frente a iteración por índices y la trampa del \textit{off-by-one}. Mutación involuntaria de objetos dentro de funciones (\textit{preview} a S9, donde se formaliza el \textit{scoping}). El criterio práctico del curso: \orange{si se puede vectorizar, se vectoriza; si no se puede, se escribe un \textit{loop} limpio}.
\end{itemize}

\medskip
\noindent\textbf{Objetivos de aprendizaje}
\begin{enumerate}[leftmargin=*, itemsep=1pt]
    \item Escribir condicionales (\texttt{if}/\texttt{else}/\texttt{switch}) y \textit{loops} (\texttt{for}/\texttt{while}) en R con sintaxis correcta.
    \item Identificar cuándo una operación se puede vectorizar y reescribirla sin \textit{loop} cuando sea posible.
    \item Pre-asignar la salida de un \textit{loop} y elegir adecuadamente entre iteración por valores e iteración por índices.
    \item Reconocer anti-patrones de iteración comunes y refactorizarlos hacia código vectorizado o, en su defecto, hacia \textit{loops} bien construidos.
\end{enumerate}

\medskip
\noindent\textbf{Lecturas}
\begin{itemize}[leftmargin=*, itemsep=1pt]
    \item \textit{Advanced R} (2.\textsuperscript{a} ed.), Cap.\ 5 \textit{Control flow} (completo) \cite{wickham2019advr}.
\end{itemize}

\semana{9}{Funciones: diseño, scoping y debugging}

Esta semana es la transición del estudiante de \textit{usuario} a \textit{autor} de código. Las funciones son la unidad fundamental de reuso y la pieza que permite consolidar el proceso ETL. Cubrimos la anatomía de una función, el modelo de \textit{scoping} con el que R resuelve nombres, las herramientas para producir y capturar errores, las técnicas para diagnosticarlos cuando ocurren, y el operador \textit{embrace} \texttt{\{\{~\}\}} ---que cierra la inquietud sembrada en la Semana 5 sobre cómo escribir funciones que interactúan con \texttt{dplyr}.

\textbf{Subtemas:}

\begin{itemize}[leftmargin=*, itemsep=2pt]
    \item \textbf{Anatomía y diseño de una función.} Sintaxis (\texttt{function(args) \{ body \}}), argumentos con valores por defecto, argumentos por nombre frente a posicionales, valor de retorno (último valor evaluado o \texttt{return()} explícito). \blue{Regla de las tres veces}: si se copia y pega tres veces el mismo patrón, conviértase en función. Convención del curso: una función tiene un solo propósito y un nombre que lo expresa.
    \item \textbf{\textit{Scoping} léxico.} Cómo R busca nombres en un orden bien definido: ambiente local $\to$ padre $\to$ global. Por qué la función no contamina el ambiente global por defecto. Variables libres y la idea de \textit{closure}. Distinción entre \blue{funciones puras} (entrada $\to$ salida, sin efectos colaterales) y funciones con \textit{side effects}. Convención del curso: preferir funciones puras siempre que sea posible.
    \item \textbf{Manejo de errores.} \texttt{stop()} para errores fatales; \texttt{warning()} para advertencias; \texttt{message()} para información. Validación temprana de argumentos al entrar a la función (\texttt{stopifnot()} o validaciones manuales con mensajes claros). Captura de errores con \texttt{tryCatch()} para flujos que deben continuar a pesar de fallas ---típico en ETL: si un trimestre falla, seguir con los siguientes y reportar al final.
    \item \textbf{Debugging.} \texttt{traceback()} para inspeccionar el \textit{call stack} tras un error. \texttt{browser()} como punto de pausa interactiva en cualquier punto del código (equivalente a un \textit{breakpoint}). \texttt{debug()}/\texttt{undebug()} para activar el modo paso a paso sobre una función completa. Estrategia general: leer el error y el \texttt{traceback} primero; \textit{print debugging} con \texttt{print()} o \texttt{cat()} segundo; debugger interactivo solo cuando sea necesario.
    \item \textbf{Funciones que reciben nombres de columnas: NSE en \texttt{dplyr}.} Por qué \texttt{mi\_funcion(df, edad)} falla cuando dentro hace \texttt{filter(df, edad > 18)}. El operador \blue{\textit{embrace}} \texttt{\{\{~var~\}\}} para reenviar un argumento al ámbito de evaluación de \texttt{dplyr}. Operador \texttt{:=} para asignar nombres dinámicos en \texttt{mutate()}. Aplicación: funciones reutilizables sobre cualquier columna del dataset, listas para incorporarse al \texttt{code/funciones.R} del proyecto integrador.
\end{itemize}

\medskip
\noindent\textbf{Objetivos de aprendizaje}
\begin{enumerate}[leftmargin=*, itemsep=1pt]
    \item Diseñar funciones de un solo propósito con argumentos bien definidos y valor de retorno explícito.
    \item Explicar cómo R resuelve los nombres dentro de una función y anticipar problemas de \textit{scoping}.
    \item Validar argumentos y producir errores informativos; capturar fallas en flujos secuenciales con \texttt{tryCatch()}.
    \item Diagnosticar errores con \texttt{traceback()} y \texttt{browser()}, distinguiendo cuándo conviene cada técnica.
    \item Escribir funciones que reciban nombres de columnas como argumentos usando el operador \texttt{\{\{~\}\}}.
\end{enumerate}

\medskip
\noindent\textbf{Lecturas}
\begin{itemize}[leftmargin=*, itemsep=1pt]
    \item \textit{Advanced R} (2.\textsuperscript{a} ed.), Cap.\ 6 \textit{Functions} (completo) \cite{wickham2019advr}.
    \item \textit{Advanced R} (2.\textsuperscript{a} ed.), Cap.\ 8 \textit{Conditions} ---§8.2 \textit{Signalling conditions} y §8.4 \textit{Handling conditions} \cite{wickham2019advr}.
    \item \textit{Advanced R} (2.\textsuperscript{a} ed.), Cap.\ 22 \textit{Debugging} (completo) \cite{wickham2019advr}.
    \item \textit{R for Data Science} (2.\textsuperscript{a} ed.), Cap.\ 25 \textit{Functions}, §25.3 \textit{Data frame functions} ---uso del operador \textit{embrace} con \texttt{dplyr} \cite{wickham2023r4ds}.
\end{itemize}

\semana{10}{Programación funcional: \texttt{purrr} y \texttt{apply}}

Esta semana cierra la Parte III del curso. Cubre la \blue{programación funcional} como paradigma: tratar funciones como objetos de datos para iterar sobre colecciones sin escribir \textit{loops}. Esta semana es continuación natural de S8 (vectorización) y S9 (funciones): donde S8 enseña a evitar \textit{loops} sobre operaciones \textit{built-in}, S10 enseña a evitar \textit{loops} sobre funciones propias.

\textbf{Subtemas:}

\begin{itemize}[leftmargin=*, itemsep=2pt]
    \item \textbf{Paradigma funcional.} Funciones como objetos de primera clase (asignables, pasables como argumentos, retornables). Por qué iterar con \textit{functionals} es preferible a un \textit{for}: predictibilidad, composición y código declarativo. Conexión con S8: si la operación es \textit{built-in}, vectoriza; si es una función propia, usa un \textit{functional}.
    \item \textbf{La familia \texttt{apply} de \textit{base R}.} \texttt{apply()} (sobre filas o columnas de una matriz), \texttt{lapply()} (lista $\to$ lista), \texttt{sapply()} (simplifica si puede ---atención: tipo de retorno impredecible), \texttt{vapply()} (con tipo declarado), \texttt{mapply()} (varios argumentos), \texttt{tapply()} (\textit{split-apply}). Limitaciones: API inconsistente y tipo de retorno no garantizado.
    \item \textbf{\texttt{purrr}: sistema unificado más moderno.} El patrón \texttt{map\_*()} con sufijo de tipo: \texttt{map()} (lista), \texttt{map\_dbl()}, \texttt{map\_chr()}, \texttt{map\_lgl()}, \texttt{map\_dfr()}/\texttt{map\_dfc()} (\textit{data frame}). Tipo de retorno predecible por construcción. Variantes para varias entradas: \texttt{map2()}, \texttt{pmap()}. Variantes para \textit{side effects}: \texttt{walk()}, \texttt{walk2()}. Funciones anónimas con \texttt{\textbackslash(x)} (R 4.1+) o \texttt{\textasciitilde{} .x} (estilo \textit{purrr}).
    \item \textbf{Reducción y predicados.} \texttt{reduce()} para colapsar una lista a un valor (acumular). \texttt{accumulate()} para mantener pasos intermedios. Predicados sobre colecciones: \texttt{keep()}, \texttt{discard()}, \texttt{every()}, \texttt{some()}, \texttt{none()}, \texttt{detect()}.
    \item \textbf{\texttt{across()}: iteración dentro de \texttt{dplyr}.} Aplicar la misma función a varias columnas dentro de \texttt{mutate()} o \texttt{summarize()}. \texttt{if\_any()} y \texttt{if\_all()} para condiciones agregadas dentro de \texttt{filter()}. Patrones aplicados al ETL: leer todos los CSVs de una carpeta, ajustar el mismo modelo a múltiples grupos, guardar N gráficas.
    \item \textbf{Filtros dinámicos en \texttt{dplyr}: \texttt{expr()} y \texttt{!!}.} Continuación natural del operador \textit{embrace} \texttt{\{\{~\}\}} introducido en S9. Cuando lo que varía no es un nombre de columna sino la \blue{expresión completa} de filtrado (e.g.\ un \texttt{shiny::input\$filtro} construido como cadena, o una lista de condiciones generadas en \texttt{map}), el camino es construir la expresión con \texttt{rlang::expr()} y \texttt{rlang::parse\_expr()}, y luego inyectarla con el operador \textit{unquote} \texttt{!!}. Sintaxis básica: \texttt{e <- expr(edad > 30 \& sexo == "M")}; \texttt{filter(datos, !!e)}. Para listas de condiciones: \texttt{filter(datos, !!!lista\_de\_exprs)} con el operador \textit{splice} \texttt{!!!}. Casos típicos: filtros parametrizados que vienen de la UI de un Shiny (preview a S15), reglas de validación que viven en un archivo de configuración, generación programática de subconjuntos para iterar con \texttt{map}.
    \item \textbf{Paralelización con \texttt{furrr} y \texttt{future}.} Una de las virtudes operativas de los \textit{functionals}: paralelizar la iteración es \blue{un cambio de prefijo}. \texttt{furrr::future\_map\_*()} reemplaza a \texttt{purrr::map\_*()} con firmas idénticas; basta declarar \texttt{future::plan(multisession, workers = N)} al inicio para distribuir las iteraciones entre N procesos. \orange{Limitación natural de R:} el lenguaje es \textit{single-threaded}; la paralelización ocurre vía procesos separados, no hilos ---no hay memoria compartida. Los objetos globales se serializan y copian a cada \textit{worker}, de modo que con N \textit{workers} y un dataset de M GB se requieren aproximadamente N $\times$ M GB de RAM. La paralelización paga cuando la tarea es \textit{CPU-bound} y suficientemente grande para amortizar el costo de \textit{spawn} y serialización; iteraciones triviales resultan más lentas en paralelo.
\end{itemize}

\medskip
\noindent\textbf{Objetivos de aprendizaje}
\begin{enumerate}[leftmargin=*, itemsep=1pt]
    \item Reconocer cuándo una iteración corresponde a un patrón funcional (mapeo, reducción, filtrado por predicado) y elegir el \textit{functional} adecuado.
    \item Usar la familia \texttt{map\_*()} de \texttt{purrr} con tipo de retorno predecible, incluyendo variantes para varias entradas.
    \item Sustituir \textit{loops} explícitos sobre archivos, columnas o grupos por \textit{functionals}, distinguiendo cuándo el \textit{loop} sigue siendo la mejor opción.
    \item Aplicar \texttt{across()} para vectorizar transformaciones sobre múltiples columnas dentro de \texttt{dplyr}.
    \item Reconocer las diferencias computacionales entre un \textit{loop} \texttt{for} y los \textit{functionals} de \texttt{purrr} (estabilidad de tipo, manejo de memoria, ámbito de variables, vía a paralelización), y usar este criterio para elegir entre uno u otro.
    \item Paralelizar una iteración con \texttt{furrr}/\texttt{future} cuando se justifica, anticipando las limitaciones de memoria que impone el modelo multi-proceso de R.
\end{enumerate}

\medskip
\noindent\textbf{Lecturas}
\begin{itemize}[leftmargin=*, itemsep=1pt]
    \item \textit{Advanced R} (2.\textsuperscript{a} ed.), Cap.\ 9 \textit{Functionals} (completo) \cite{wickham2019advr}.
    \item \textit{R for Data Science} (2.\textsuperscript{a} ed.), Cap.\ 26 \textit{Iteration} (completo) \cite{wickham2023r4ds}.
\end{itemize}

\medskip
\noindent\textbf{Referencia rápida:} \cheatsheet{purrr}{Apply Functions with \texttt{purrr}} (Posit) ---familia \texttt{map\_*}, variantes y patrones de iteración funcional en una página.

% ------------------------------------------------------------------------------
\subsection*{\color{main}Parte IV --- Análisis}
% ------------------------------------------------------------------------------

\semana{11}{Modelado estadístico y series de tiempo}

Esta semana abre la Parte IV del curso. El enfoque es \blue{estrictamente programático}: cómo se especifican fórmulas, cómo se ajustan modelos, cómo se extraen y reportan resultados. La teoría estadística detrás de cada modelo se asume conocida ---el curso de Econometría es prerrequisito conceptual para esta semana. La meta práctica es que el estudiante pueda mover una pregunta de su proyecto integrador (¿qué predice X?, ¿hay diferencia entre grupos?, ¿cómo evoluciona Y en el tiempo?) hasta una tabla \textit{publication-ready} de resultados.

\textbf{Subtemas:}

\begin{itemize}[leftmargin=*, itemsep=2pt]
    \item \textbf{Fórmulas y \texttt{lm()}.} Sintaxis de fórmulas: operadores (\texttt{+}, \texttt{:}, \texttt{*}, \texttt{\textasciicircum}, \texttt{I()}), control del intercepto (\texttt{-1}/\texttt{+0}). Anatomía del objeto \texttt{lm}: \texttt{summary()}, \texttt{coef()}, \texttt{confint()}, \texttt{predict()}, \texttt{residuals()}, \texttt{fitted()}. Construcción incremental con \texttt{update()}. Diagnósticos básicos: \texttt{plot.lm()} (los cuatro paneles canónicos) y \texttt{vif()} para multicolinealidad.
    \item \textbf{Categóricos en modelos.} \textbf{Como predictores:} \textit{factors} convertidos automáticamente a \textit{dummies}, manejo del \textit{reference level} con \texttt{relevel()} o reordenando el \textit{factor}, interpretación de coeficientes. \textbf{Como variable dependiente:} \texttt{glm(family = binomial)} para modelos logit/probit, \texttt{glm(family = poisson)} para conteos. Interpretación operativa de coeficientes (\textit{odds-ratio}, exponenciación) sin entrar a teoría estadística.
    \item \textbf{Datos panel y errores estándar robustos.} \texttt{pdata.frame()} con identificador y dimensión temporal. Estimadores: \texttt{model = 'within'} (FE), \texttt{'random'} (RE), \texttt{'between'}, \texttt{'pooling'}. Tests de Hausman y Breusch-Pagan vía \texttt{phtest()} y \texttt{plmtest()} (uso programático, sin demostraciones). Errores estándar robustos a heteroscedasticidad con \texttt{sandwich::vcovHC()} y reporte vía \texttt{lmtest::coeftest()}, aplicables tanto a \texttt{lm} como a \texttt{plm}.
    \item \textbf{Fechas y series de tiempo (nivel mínimo).} Fechas con \texttt{lubridate}: \textit{parsers} (\texttt{ymd()}, \texttt{dmy()}, \texttt{mdy()}), extractores (\texttt{year}, \texttt{month}, \texttt{quarter}, \texttt{wday}), aritmética (\texttt{days()}, \texttt{months()}, \texttt{years()}) y secuencias temporales. Clases básicas: \texttt{ts} (regular), \texttt{zoo} (irregular), \texttt{xts} (heredera moderna); conversión entre \texttt{tibble} y series. Operaciones de diagnóstico: \texttt{lag()}, \texttt{diff()}, ACF y PACF visuales. 
    \item \textbf{Reportes con \texttt{broom} y \texttt{modelsummary}.} \texttt{broom}: \texttt{tidy()} para coeficientes en \textit{tibble}, \texttt{glance()} para \textit{fit statistics} en una fila, \texttt{augment()} para residuales y \textit{fitted} añadidos al \textit{data}. La misma API funciona para \texttt{lm}, \texttt{glm}, \texttt{plm} y muchos otros. \texttt{modelsummary}: comparación de varios modelos lado a lado en una tabla, customización (etiquetas, estrellas, decimales), \textit{output} a Word, HTML, Markdown y LaTeX. Comparación con \texttt{stargazer}, herramienta tradicional en economía académica.
\end{itemize}

\medskip
\noindent\textbf{Objetivos de aprendizaje}
\begin{enumerate}[leftmargin=*, itemsep=1pt]
    \item Especificar modelos lineales con la sintaxis de fórmulas de R, manipular el objeto \texttt{lm} y producir diagnósticos básicos.
    \item Incorporar variables categóricas en modelos como predictores y como variable dependiente (vía GLM), interpretando los coeficientes en términos operativos.
    \item Estimar modelos de datos panel con \texttt{plm} y reportar errores estándar robustos a heteroscedasticidad con \texttt{sandwich} + \texttt{lmtest}.
    \item Manejar fechas con \texttt{lubridate} y las clases básicas de series de tiempo (\texttt{ts}, \texttt{zoo}), aplicando operaciones de diagnóstico (\textit{lag}, diferencias, ACF/PACF).
    \item Producir tablas de modelos \textit{publication-ready} con \texttt{broom} y \texttt{modelsummary}, comparando varios modelos lado a lado.
\end{enumerate}

\medskip
\noindent\textbf{Lecturas}
\begin{itemize}[leftmargin=*, itemsep=1pt]
    \item \textit{An Introduction to R} (manual oficial), Cap.\ 11 \textit{Statistical models in R} \cite{rcoreteam-intro}.
    \item \textit{Panel Data Econometrics in R: The plm Package} (Croissant \& Millo) ---viñeta principal del paquete \cite{croissant-plm}.
    \item Documentación oficial de \texttt{broom} \cite{robinson-broom} y \texttt{modelsummary} \cite{arelbundock-modelsummary} ---secciones \textit{Getting started}.
\end{itemize}

\medskip
\noindent\textbf{Referencia rápida:} \cheatsheet{lubridate}{Dates and times with \texttt{lubridate}} (Posit) ---\textit{parsers}, extractores y aritmética temporal en una página.

\semana{12}{Datos a escala: SQL, \texttt{DBI} y \texttt{dbplyr}}

Esta semana cierra la Parte IV del curso. Cubre las herramientas para conectar R con bases de datos relacionales: SQL como lenguaje, \texttt{DBI} como API de conexión y \texttt{dbplyr} como puente con la sintaxis de \texttt{dplyr}. El objetivo es conectar el \textit{Load} del ETL a una base local y aplicar análisis estadísticos sobre esos datos, de modo que el estudiante pueda mover datos entre archivos, memoria y BD, y elegir con criterio cuál es la mejor representación para cada momento del proyecto.

\textbf{Subtemas:}

\begin{itemize}[leftmargin=*, itemsep=2pt]
    \item \textbf{Motivación y tipos de bases de datos.} Por qué una BD: datos que no caben en memoria, persistencia entre ejecuciones, varios procesos accediendo a los mismos datos, fuentes externas que se exponen vía SQL (CONEVAL, catálogos del INEGI, Banxico). Comparación práctica: \blue{SQLite} (archivo, \textit{zero-config}, el \textit{default} del curso); \texttt{PostgreSQL} y \texttt{MySQL} (servidores, lo común en producción real); \blue{duckdb} como alternativa moderna ---\textit{zero-config} como SQLite pero columnar y optimizada para \textit{analytics}, con \textit{drivers} excelentes en R; particularmente atractiva para encuestas con varios millones de filas. Mención a \texttt{arrow}/\textit{parquet} como alternativa \textit{file-based} para \textit{analytics} a escala (no es BD pero compite en el nicho).
    \item \textbf{SQL básico.} \texttt{SELECT}, \texttt{FROM}, \texttt{WHERE}, \texttt{ORDER BY}, \texttt{LIMIT}. \texttt{JOIN} (\texttt{INNER}, \texttt{LEFT}, \texttt{RIGHT}, \texttt{FULL}). \texttt{GROUP BY} con funciones agregadoras (\texttt{COUNT}, \texttt{SUM}, \texttt{AVG}, \texttt{MIN}, \texttt{MAX}). Operadores: \texttt{LIKE}, \texttt{IN}, \texttt{BETWEEN}, \texttt{IS NULL}. \textit{Common Table Expressions} con \texttt{WITH} para queries legibles que se construyen en pasos. Mención a \textit{subqueries}. Statements de creación: \texttt{CREATE TABLE} e \texttt{INSERT} ---solo lo necesario para el \textit{Load} del proyecto.
    \item \textbf{\texttt{DBI}: la API de conexión.} \texttt{dbConnect()}/\texttt{dbDisconnect()} para abrir y cerrar conexiones; conveniencia de envolver en \texttt{tryCatch()} para garantizar el cierre. Operaciones de tabla: \texttt{dbWriteTable()}, \texttt{dbReadTable()}, \texttt{dbListTables()}, \texttt{dbListFields()}. Ejecución de queries: \texttt{dbGetQuery()} cuando se espera resultado tabular; \texttt{dbExecute()} cuando no (creación, inserción, actualización). Parametrización segura: \texttt{glue\_sql()} para interpolar valores en queries sin riesgo de \textit{SQL injection}.
    \item \textbf{\texttt{dbplyr}: \texttt{dplyr} sobre SQL.} El concepto: \texttt{tbl(con, 'tabla')} crea un \textit{remote tibble} sobre el cual los verbos de \texttt{dplyr} se traducen automáticamente a SQL y se ejecutan en el \textit{server}. \texttt{show\_query()} para inspeccionar el SQL generado ---herramienta de aprendizaje invaluable. \texttt{collect()} para materializar el resultado en memoria como \textit{tibble} ordinario. \blue{Patrón clave}: hacer \texttt{filter}/\texttt{group\_by}/\texttt{summarize} en el \textit{server} (que es bueno y rápido en eso) y solo \texttt{collect()} al final, cuando el resultado ya es chico. Limitación honesta: no todas las funciones de R se traducen ---funciones \textit{custom} no van.
     \textit{Preview} de S15: el Shiny consumirá esta misma BD como fuente.
\end{itemize}

\medskip
\noindent\textbf{Objetivos de aprendizaje}
\begin{enumerate}[leftmargin=*, itemsep=1pt]
    \item Reconocer cuándo un proyecto de datos justifica una BD relacional y elegir entre las opciones disponibles (SQLite, duckdb, PostgreSQL, MySQL) con criterio.
    \item Escribir queries SQL básico que cubran el grueso de las necesidades de análisis (\texttt{SELECT}/\texttt{JOIN}/\texttt{GROUP BY}/agregadoras/\texttt{WITH}).
    \item Conectar R a una BD vía \texttt{DBI}, escribir y leer tablas, y ejecutar queries parametrizadas de forma segura.
    \item Usar \texttt{dbplyr} para combinar la sintaxis de \texttt{dplyr} con la ejecución en el \textit{server} de BD, e inspeccionar el SQL generado.
\end{enumerate}

\medskip
\noindent\textbf{Lecturas}
\begin{itemize}[leftmargin=*, itemsep=1pt]
    \item \textit{R for Data Science} (2.\textsuperscript{a} ed.), Cap.\ 21 \textit{Databases} (completo) \cite{wickham2023r4ds}.
    \item Documentación oficial de \texttt{DBI} ---sección \textit{Get started}--- y \texttt{dbplyr} ---viñetas \textit{Introduction to dbplyr} y \textit{Verb translation}.
\end{itemize}

% ------------------------------------------------------------------------------
\subsection*{\color{main}Parte V --- Productos}
% ------------------------------------------------------------------------------

\semana{13}{Visualización con \texttt{ggplot2}}

Esta semana abre la Parte V (Productos). Cubre el detalle de \texttt{ggplot2} no como un catálogo de funciones para hacer gráficas, sino como la \blue{gramática de gráficas} ---la idea de que una visualización es la composición de capas independientes, cada una con responsabilidades acotadas. La meta pedagógica es que el estudiante pueda razonar por capa: si una gráfica no se ve como quiere, identifica qué capa controla qué y modifica solo esa. Esta forma de pensar es lo que distingue al usuario que copia plantillas del que produce visualizaciones a medida.

\textbf{Subtemas:}

\begin{itemize}[leftmargin=*, itemsep=2pt]
    \item \textbf{Gramática de gráficas y arquitectura por capas.} Origen del enfoque (\textit{Grammar of Graphics} de Wilkinson). Una gráfica como composición de \blue{capas independientes}, cada una controlable por separado:
    \begin{itemize}[leftmargin=*, itemsep=1pt, label=$\circ$]
        \item \textbf{Data} ---el \textit{tibble} que se grafica.
        \item \textbf{Aesthetics} (\texttt{aes}) ---mapeo variable a propiedad visual (\texttt{x}, \texttt{y}, \texttt{color}, \texttt{fill}, \texttt{size}, \texttt{shape}, \texttt{alpha}, \texttt{linetype}).
        \item \textbf{Geoms} (\texttt{geom\_*}) ---la marca visual: puntos, líneas, barras, áreas.
        \item \textbf{Stats} (\texttt{stat\_*}) ---transformaciones estadísticas implícitas (\textit{count}, \textit{smooth}, \textit{bin}).
        \item \textbf{Position adjustments} ---\texttt{stack}, \texttt{dodge}, \texttt{jitter}, \texttt{fill}.
        \item \textbf{Scales} (\texttt{scale\_*}) ---cómo los \textit{aesthetics} se traducen a valores visuales.
        \item \textbf{Coordinate systems} (\texttt{coord\_*}) ---cartesiano, polar, transformaciones de coordenadas.
        \item \textbf{Facets} (\texttt{facet\_*}) ---particionar la gráfica en sub-paneles.
        \item \textbf{Theme} ---todo lo no-dato: tipografía, fondos, ejes, leyendas.
    \end{itemize}
    \item \textbf{Geometrías, estéticos y la distinción \textit{setting/mapping}.} \textit{Geoms} de uso frecuente: \texttt{geom\_point}, \texttt{geom\_line}, \texttt{geom\_bar}/\texttt{geom\_col} (la distinción importa: \texttt{bar} cuenta, \texttt{col} grafica el valor), \texttt{geom\_histogram}, \texttt{geom\_boxplot}, \texttt{geom\_density}, \texttt{geom\_smooth}, \texttt{geom\_text}/\texttt{geom\_label}, \texttt{geom\_ribbon}, \texttt{geom\_area}, \texttt{geom\_jitter}. \blue{Distinción crítica}: \texttt{color = grupo} dentro de \texttt{aes()} es un \textit{mapping} (la variable controla el color); \texttt{color = 'red'} fuera de \texttt{aes()} es un \textit{setting} (color fijo). Confundirlos es la fuente \#1 de bugs en \texttt{ggplot2}.
    \item \textbf{Scales: el control fino sobre cómo se ve cada estético.} \texttt{scale\_x\_continuous}, \texttt{scale\_x\_discrete}, \texttt{scale\_x\_log10}, \texttt{scale\_x\_date}; análogos para \texttt{y}, \texttt{color}, \texttt{fill}, \texttt{size}. Control de \texttt{limits}, \texttt{breaks}, \texttt{labels}, transformaciones. \texttt{scale\_*\_manual()} con vectores nombrados para control total. \blue{El paquete \texttt{scales}} como complemento esencial: formatters como \texttt{label\_dollar()}, \texttt{label\_percent()}, \texttt{label\_comma()}, \texttt{label\_number\_si()} convierten ejes amateur en ejes profesionales. Diferencia importante entre \texttt{coord\_cartesian(ylim = ...)} (zoom) y \texttt{scale\_y\_continuous(limits = ...)} (descarta datos).
    \item \textbf{Facets, coordenadas y themes.} \texttt{facet\_wrap} vs.\ \texttt{facet\_grid}; cuándo usar \texttt{scales = 'free'} y cuándo no. \texttt{coord\_flip()} para barras horizontales, \texttt{coord\_polar()} para circulares, \texttt{coord\_cartesian()} para zoom sin reajuste. \textbf{Themes:} predefinidos (\texttt{theme\_minimal}, \texttt{theme\_bw}, \texttt{theme\_classic}, \texttt{theme\_void}) y customización fina con \texttt{theme()} ---ajuste de \texttt{axis.text}, \texttt{panel.grid}, \texttt{legend.position}, \texttt{plot.title}, etc. \texttt{theme\_set()} para fijar tema global del proyecto.
    \item \textbf{Ecosistema complementario y exportación.} \textbf{Composición y anotación:} \texttt{patchwork} (combinar varias gráficas con \texttt{+}, \texttt{/} y \texttt{|}), \texttt{ggrepel} (\texttt{geom\_text\_repel} para labels que no se sobreponen), \texttt{ggtext} (\textit{markdown}/HTML en títulos y labels), \texttt{gghighlight} (destacar grupos en gráficas multi-grupo). \textbf{Color:} \texttt{ggthemes} (\textit{themes} estilo \textit{Economist}, \textit{FiveThirtyEight}, \textit{WSJ}), \texttt{paletteer} (acceso unificado a cientos de paletas). \textit{Mención somera}: \texttt{plotly} para gráficas interactivas, \texttt{gganimate} para animaciones. \textbf{Exportación con \texttt{ggsave()}:} formatos PNG (\textit{raster}, ideal para web y Word), JPG (\textit{raster}, calidad inferior, evitar salvo necesidad) y SVG (vector, escalable sin pérdida). Argumentos clave: \texttt{width}, \texttt{height}, \texttt{units} (cm/in/mm/px), \texttt{dpi}, \texttt{bg = 'white'} para forzar fondo blanco.
\end{itemize}

\medskip
\noindent\textbf{Objetivos de aprendizaje}
\begin{enumerate}[leftmargin=*, itemsep=1pt]
    \item Explicar la arquitectura por capas de \texttt{ggplot2} y nombrar qué controla cada una.
    \item Construir una gráfica eligiendo la geometría adecuada y distinguiendo entre \textit{mapping} (\texttt{aes()}) y \textit{setting}.
    \item Customizar \textit{scales} (transformaciones, \textit{breaks}, \textit{labels}, formatters de \texttt{scales}) para producir ejes profesionales.
    \item Aplicar \textit{themes}, particionar con \textit{facets} y elegir la librería complementaria adecuada para cada necesidad (composición, labels sin solape, paletas, etc.).
    \item Exportar gráficas con dimensiones y formato apropiados al medio de destino (PNG, JPG, SVG).
\end{enumerate}

\medskip
\noindent\textbf{Lecturas}
\begin{itemize}[leftmargin=*, itemsep=1pt]
    \item \textit{R for Data Science} (2.\textsuperscript{a} ed.), Cap.\ 9 \textit{Layers} (completo) \cite{wickham2023r4ds}.
    \item \textit{R for Data Science} (2.\textsuperscript{a} ed.), Cap.\ 11 \textit{Communication} (completo; nota: el capítulo cubre \textit{labels}, \textit{scales}, \textit{themes} y \textit{layout} con \texttt{patchwork}, pero \textbf{no} cubre \texttt{ggsave}/exportación ---esa parte se cubre directamente en clase y con la documentación del paquete) \cite{wickham2023r4ds}.
    \item \textit{ggplot2: Elegant Graphics for Data Analysis} (Wickham) ---referencia comprehensiva, consultar capítulos relevantes según necesidad \cite{wickham2016ggplot2}.
\end{itemize}

\medskip
\noindent\textbf{Referencia rápida:} \cheatsheet{data-visualization}{Data visualization with \texttt{ggplot2}} (Posit) ---arquitectura por capas, geometrías comunes, \textit{scales} y \textit{themes} en una página.

\semana{14}{Automatización con \texttt{officer} y \texttt{googledrive}}

Esta semana se enfoca en convertir los outputs del análisis (gráficas, tablas, valores clave) en \blue{presentaciones automatizadas} y en establecer un canal de \blue{gestión de archivos compartidos} con Google Drive. La meta práctica es desplazar el patrón habitual ---copiar y pegar gráficas en PowerPoint manualmente cada vez que cambian los datos--- por un flujo donde la presentación se regenera con un \textit{script}, partiendo de un \textit{template} estilizado, y se publica al equipo automáticamente. El enfoque del curso privilegia PPT sobre Word como medio de entrega, en línea con la práctica más frecuente en análisis aplicado.

\textbf{Subtemas:}

\begin{itemize}[leftmargin=*, itemsep=2pt]
    \item \textbf{Motivación y arquitectura del flujo automatizado.} El problema: tras cada actualización de datos, alguien rehace la presentación. Costo en tiempo, fragilidad, dificultad para auditar. La automatización resuelve esto con la fórmula \texttt{template + datos + código = presentación}. Patrón general: el pipeline ETL produce outputs en \texttt{output/}; un \textit{script} de presentación lee un \texttt{.pptx} \textit{template} pre-estilizado e inyecta los valores actualizados; el resultado se publica vía Drive.
    \item \textbf{\texttt{officer} para PowerPoint.} \texttt{read\_pptx()} para cargar un \textit{template} (recomendado siempre sobre construir desde cero). Agregar diapositivas con \texttt{add\_slide()} eligiendo \textit{layouts} del \textit{master} (Título, Contenido, Sección, etc.). Insertar contenido en \textit{placeholders} con nombre vía \texttt{ph\_with()}: texto, tablas, gráficas \texttt{ggplot2}, imágenes. Iteración sobre múltiples diapositivas para presentaciones con N secciones (un cuadro por estado, por trimestre, etc.). Exportación final con \texttt{print(doc, target = ...)}. Mención: \texttt{officer} también genera Word, aunque está fuera del énfasis del curso.
    \item \textbf{Estilización de componentes.} Donde la presentación pasa de funcional a profesional. Texto enriquecido: \texttt{fp\_text()} (color, tamaño, negrita, cursiva), \texttt{fp\_par()} (alineación, espaciado), \texttt{fpar()} y \texttt{block\_list()} para composiciones inline (mezclar negrita + color + tamaño en el mismo párrafo). Tablas con \texttt{flextable}: themes (\texttt{theme\_vanilla}, \texttt{theme\_box}, custom), formato numérico por columna (\texttt{colformat\_double}, \texttt{colformat\_pct}), \textit{merging} de celdas, \textit{headers} compuestos, colores por celda. Inserción de gráficas con dimensiones controladas (\texttt{ph\_with(value = gg, location = ph\_location(width, height))}). La consigna pedagógica: \blue{todo el estilo vive en el template y en código de \texttt{officer}, no se ``arregla'' en PowerPoint a mano} (eso destruye la reproducibilidad).
    \item \textbf{\texttt{googledrive}: autenticación y gestión de archivos.} \textbf{Autenticación:} OAuth interactivo para desarrollo (\texttt{drive\_auth()}); \textit{Service Account} con JSON para automatización (\texttt{drive\_auth(path = "credenciales.json")}). \textbf{Gestión de archivos como solución \textit{low-code} de automatización:} \texttt{drive\_upload()} para publicar reportes generados; \texttt{drive\_download()} para descargar insumos compartidos; \texttt{drive\_ls()} para navegar carpetas; \texttt{drive\_get()} para localizar archivos por nombre o ID; \texttt{drive\_share()} para gestionar permisos; \texttt{drive\_mv()} y \texttt{drive\_trash()} para mover y archivar. \blue{Patrón típico}: una carpeta de Drive funciona como \textit{buzón} ---el equipo deposita insumos ahí, el \textit{script} los descarga automáticamente; el reporte generado se sube a otra carpeta donde el equipo lo consume sin contacto con el código. Esto reemplaza muchos flujos manuales con muy poca complejidad técnica.
    \item \textbf{\texttt{gmailr}: correo como canal de automatización.} Autenticación análoga a \texttt{googledrive} (OAuth interactivo en desarrollo; \textit{Service Account} para automatización) sobre la API de Gmail. \textbf{Envío programático:} \texttt{gm\_mime()} para componer un correo (asunto, cuerpo en HTML, destinatarios, adjuntos), \texttt{gm\_send\_message()} para enviarlo; útil para distribuir el reporte generado a una lista de \textit{stakeholders} cada vez que corre el pipeline. \textbf{Lectura como buzón de entrada:} \texttt{gm\_messages()} para listar correos según filtros tipo Gmail (\texttt{from:}, \texttt{subject:}, \texttt{has:attachment}), \texttt{gm\_message()} para obtener uno, \texttt{gm\_attachments()} y \texttt{gm\_save\_attachments()} para extraer adjuntos al disco. \blue{Patrón \textit{low-code}}: el \textit{script} lee la bandeja con un filtro específico (e.g.\ \texttt{from:proveedor@empresa.com subject:'datos mensuales'}), descarga los adjuntos, los procesa y manda el reporte por correo al equipo. Reemplaza intercambios manuales de archivos por algo que solo requiere que el equipo siga usando su correo electrónico habitual.
\end{itemize}

\medskip
\noindent\textbf{Objetivos de aprendizaje}
\begin{enumerate}[leftmargin=*, itemsep=1pt]
    \item Justificar el patrón \texttt{template + datos + código = presentación} frente al flujo manual de copiar y pegar.
    \item Generar diapositivas de PowerPoint programáticamente a partir de un \textit{template}, insertando texto, tablas y gráficas en \textit{placeholders} con nombre.
    \item Estilizar componentes de una presentación (texto enriquecido, tablas con \texttt{flextable}, imágenes dimensionadas) sin recurrir a edición manual en PowerPoint.
    \item Operar Google Drive desde R para descarga, carga y gestión de archivos, distinguiendo OAuth interactivo de \textit{Service Account}.
    \item Usar \texttt{gmailr} para enviar reportes generados y consumir una bandeja de Gmail como fuente de adjuntos.
    \item Diseñar canales de automatización \textit{low-code} en los que carpetas de Drive y bandejas de Gmail funcionan como buzones de insumos/outputs entre el código y los estudiantes.
\end{enumerate}

\medskip
\noindent\textbf{Lecturas}
\begin{itemize}[leftmargin=*, itemsep=1pt]
    \item Documentación oficial de \texttt{officer} ---viñetas \textit{Get started with PowerPoint} y \textit{Layouts and slide manipulation} \cite{gohel-officer}.
    \item Documentación oficial de \texttt{flextable} ---secciones \textit{Quick start} y \textit{Theme functions} \cite{gohel-flextable}.
    \item Documentación oficial de \texttt{googledrive} ---viñeta \textit{Authentication} y referencia de \textit{drive\_*} \cite{bryan-googledrive}.
    \item Documentación oficial de \texttt{gmailr} ---viñetas \textit{Sending email} y \textit{Working with messages} \cite{hester-gmailr}.
\end{itemize}

\semana{15}{Aplicaciones interactivas con Shiny}

Esta semana es introducción, no dominio. Shiny es una tecnología vasta y profesionalmente seria; el objetivo del curso no es formar desarrolladores web sino \blue{dar los fundamentos suficientes para que el estudiante construya un tablero funcional} (filtros, gráficas interactivas, indicadores) que sirva como producto final del proyecto integrador. La meta es que entiendan qué hay debajo del capó ---HTML, CSS, la capa de conversión, la reactividad--- y que reconozcan las limitaciones de la herramienta. Quien quiera profundizar después tendrá las bases para hacerlo.

\textbf{Subtemas:}

\begin{itemize}[leftmargin=*, itemsep=2pt]
    \item \textbf{Anatomía de una página web: HTML y CSS (lo mínimo).} Modelo \textit{cliente-servidor}: el navegador pide, el servidor responde con HTML/CSS/JavaScript. \textbf{HTML} como árbol jerárquico de \textit{tags} (\texttt{div}, \texttt{span}, \texttt{h1}--\texttt{h6}, \texttt{p}, \texttt{a}, \texttt{img}, \texttt{ul}/\texttt{li}); atributos \texttt{id} y \texttt{class}. \textbf{CSS} como sistema de reglas para dar estilo: selectores por \textit{tag}, \texttt{class}, \texttt{id}; propiedades comunes (\texttt{color}, \texttt{font}, \texttt{background}, \texttt{margin}, \texttt{padding}); cuándo va \textit{inline} y cuándo en hoja externa. No buscamos formar expertos en web ---solo dar la intuición de qué hay debajo de cualquier aplicación Shiny.
    \item \textbf{Cómo Shiny convierte R a HTML/CSS/JS.} Funciones de R que producen HTML: \texttt{tags\$div()}, \texttt{tags\$span()}, \texttt{HTML()}, \texttt{h1()}, \texttt{p()}, \texttt{a()}. Los \textit{inputs} y \textit{outputs} son elementos HTML especiales con atributos que el JavaScript del cliente escucha. \textbf{Bootstrap} como la librería de CSS que Shiny carga por defecto (de ahí el sistema de \textit{grid} y los estilos típicos). \textbf{\texttt{htmlwidgets}} como el puente para integrar librerías JavaScript (\texttt{highcharts}, \texttt{plotly}, \texttt{leaflet}, etc.) sin escribir JavaScript.
    \item \textbf{Anatomía mínima de una app Shiny.} Separación \textbf{UI} (qué se ve) vs. \textbf{server} (qué hace). \texttt{shinyApp(ui, server)} para definir y \texttt{runApp()} para ejecutar. \textit{Inputs} declarados en UI (\texttt{selectInput}, \texttt{numericInput}, \texttt{sliderInput}, \texttt{dateInput}, \texttt{checkboxInput}, \texttt{textInput}) accesibles en \textit{server} vía \texttt{input\$id}. \textit{Outputs} declarados en UI (\texttt{plotOutput}, \texttt{tableOutput}, \texttt{textOutput}) poblados en \textit{server} con \texttt{output\$id <- render*(\{...\})}. \blue{Reactividad} como modelo de programación: \texttt{reactive()} para valores que se recalculan al cambiar sus inputs; \texttt{observe()}/\texttt{observeEvent()} para \textit{side effects}; \texttt{reactiveVal()} para estado mutable explícito; \texttt{isolate()} para romper la reactividad cuando se necesita.
    \item \textbf{Estructura tipo tablero con \texttt{bslib}.} El paquete moderno para layouts de \textit{dashboard}. \texttt{page\_sidebar()} para la estructura clásica (filtros a la izquierda, contenido principal a la derecha); \texttt{page\_navbar()} para tableros multi-sección. \texttt{value\_box()} para indicadores numéricos clave (KPIs). \texttt{card()} y \texttt{card\_header()} para encapsular secciones. \texttt{layout\_column\_wrap()} y \texttt{layout\_columns()} para grid responsive. \textbf{Gráficas interactivas} con \texttt{highcharter} (estética profesional, gran interactividad nativa) o \texttt{plotly} (más integrado al ecosistema \textit{tidyverse} vía \texttt{ggplotly()}). Los filtros se conectan a las gráficas vía reactividad: cambia el filtro $\to$ se recalcula la gráfica.
    \item \textbf{Limitaciones y cuándo Shiny es la herramienta correcta.} Cada sesión de Shiny es un proceso R completo: N usuarios concurrentes implica N veces la memoria. La escalabilidad no es nativa; servir muchos usuarios requiere infraestructura no trivial (Posit Connect, Docker). El modelo reactivo se vuelve complejo rápido si se busca interactividad fina. \orange{Shiny brilla cuando}: el output es R-céntrico (datos, gráficas, modelos), el público es interno o limitado, el tiempo de desarrollo es la restricción, y la app es para exploración o reporte interactivo. \orange{Shiny no es la herramienta} cuando: hay que escalar a miles de usuarios concurrentes, la interactividad pedida es muy granular (entonces React/Vue), o el público no tiene relación con R. \textbf{Despliegue del curso:} \texttt{shinyapps.io} (gratuito con límite de horas/mes, suficiente para mostrar el proyecto).
\end{itemize}

\medskip
\noindent\textbf{Objetivos de aprendizaje}
\begin{enumerate}[leftmargin=*, itemsep=1pt]
    \item Explicar la arquitectura cliente-servidor de una página web y reconocer los roles de HTML, CSS y JavaScript.
    \item Identificar cómo las funciones de Shiny generan HTML/CSS/JS bajo el capó (\texttt{tags\$*}, Bootstrap, \texttt{htmlwidgets}).
    \item Construir una aplicación Shiny mínima con UI, \textit{server} e interacción reactiva.
    \item Implementar un tablero con filtros, \texttt{value\_box}es y gráficas interactivas usando \texttt{bslib} junto con \texttt{highcharter} o \texttt{plotly}.
    \item Discernir cuándo Shiny es la herramienta apropiada y cuándo conviene mirar a otras tecnologías.
\end{enumerate}

\medskip
\noindent\textbf{Lecturas}
\begin{itemize}[leftmargin=*, itemsep=1pt]
    \item \textit{Mastering Shiny} (Wickham), capítulos introductorios sobre UI básica, \textit{server} y reactividad \cite{wickham2021mastering-shiny}.
    \item Documentación oficial de \texttt{bslib} ---viñeta \textit{Dashboards} \cite{bslib-docs}.
    \item Documentación oficial de \texttt{highcharter} (\textit{Get started}) \cite{kunst-highcharter} \textbf{o} \textit{Interactive Web-Based Data Visualization with R, plotly, and shiny} (Sievert), capítulos introductorios \cite{sievert-plotly} ---según la opción de visualización elegida.
\end{itemize}

\medskip
\noindent\textbf{Referencia rápida:} \cheatsheet{shiny}{Web Applications with \texttt{shiny}} (Posit) ---estructura de la app, \textit{inputs}/\textit{outputs}, reactividad y \textit{layouts} en una página.

\semana{16}{Interfaces de modelos de lenguaje}

Esta es la última semana del curso y está concentrada en el límite de las capacidades programáticas. El tema: cómo incorporar modelos de lenguaje (LLMs) al pipeline de datos como una herramienta más, en lugar de tratarlos como un servicio externo que rompe el flujo. La meta práctica es que el estudiante pueda mover tareas que habitualmente se delegan manualmente a un \textit{chatbot} ---extraer información estructurada de un párrafo, clasificar texto, generar un resumen, normalizar nombres--- al interior del pipeline, donde son \blue{auditables, reproducibles y se ejecutan sin intervención humana}.

\textbf{Subtemas:}

\begin{itemize}[leftmargin=*, itemsep=2pt]
    \item \textbf{Motivación: LLMs integrados al pipeline.} El patrón roto que vamos a sustituir: el analista copia un párrafo no estructurado, lo pega en ChatGPT/Claude/Gemini, pide ``extrae estos campos''/``resume''/``clasifica'', copia la respuesta y la pega al dataset. Problemas: no reproducible, no escala, no auditable, rompe el flujo. La alternativa: el mismo LLM \textit{adentro} del \textit{script}, llamado como cualquier otra función. Cada corrida del pipeline ejecuta la inferencia automáticamente; los resultados quedan en el dataframe; el código documenta exactamente qué se le pidió al modelo.
    \item \textbf{Dos formas de llamar a una API de LLM.} \textbf{Forma manual con \texttt{httr2}:} construir el \textit{request} POST a \texttt{/v1/chat/completions} (o el endpoint equivalente del proveedor), incluir \textit{headers} con autenticación, armar el JSON del \textit{body}, ejecutar y parsear la respuesta. Ventajas: control total, sin dependencias adicionales, transparencia sobre qué se manda y qué llega. Desventajas: verboso, requiere conocer la estructura de cada API. \textbf{Forma amigable con \texttt{ellmer}:} el paquete de Posit que abstrae la API en un objeto \texttt{chat} con métodos consistentes; soporta múltiples proveedores (Anthropic, OpenAI, Google Gemini, Groq, Ollama local, etc.) cambiando solo la función de inicialización. Ventajas: portabilidad entre proveedores, API limpia, soporte para \textit{streaming} y para herramientas estructuradas. Desventajas: una capa adicional que oculta detalles. \blue{Recomendación del curso:} \texttt{ellmer} para uso productivo, \texttt{httr2} cuando se necesita control fino o se aprende qué hay debajo.
    \item \textbf{Salidas estructuradas: del texto al dato.} El uso más valioso de un LLM en un pipeline no es obtener texto libre, sino \blue{datos estructurados}. \texttt{ellmer} soporta esto vía \texttt{chat\$chat\_structured()} con un \textit{schema} declarado con \texttt{type\_object()}, \texttt{type\_string()}, \texttt{type\_number()}, \texttt{type\_array()}, etc. El modelo es forzado a responder en JSON que matchea el \textit{schema}, lo cual se parsea directamente a un \textit{tibble}. Esto convierte el LLM en un transformador determinado-en-forma: entrada texto $\to$ salida con campos fijos. Aplicación inmediata: extraer N campos de un párrafo (e.g.\ de un acta, una noticia, una respuesta abierta de encuesta) en una sola pasada que produce una fila lista para \texttt{bind\_rows()}.
    \item \textbf{Casos de uso típicos en proyectos de datos.} \textbf{Extracción estructurada:} párrafo $\to$ tibble con campos definidos. \textbf{Clasificación:} texto $\to$ categoría de una lista cerrada (sentimiento, tema, polaridad). \textbf{Resumen:} texto largo $\to$ resumen corto con longitud controlada. \textbf{Normalización:} texto sucio $\to$ texto canónico (nombres de municipios, instituciones, productos). \textbf{Etiquetado:} descripción $\to$ lista de \textit{tags}. \textbf{Traducción:} texto en un idioma $\to$ otro. La idea general: cualquier transformación texto-a-algo que un humano haría rápido pero tedioso para N filas, se puede automatizar con un LLM y un buen \textit{prompt}.
    \item \textbf{Panorama: extensiones del flujo.} Una vez que el estudiante entiende el patrón ``LLM como función dentro del pipeline'', el panorama de aplicaciones se abre. Solo en mención superficial: \textbf{modelos locales con \texttt{ollama}} ---correr LLMs en el propio equipo sin enviar datos a un servidor externo, útil para información sensible; \textbf{OCR local} ---extraer texto de imágenes/PDFs escaneados con paquetes como \texttt{tesseract}; \textbf{censura/redacción de datos personales} ---usar LLMs o regex con apoyo de LLM para detectar y eliminar PII (CURPs, RFCs, direcciones) antes de publicar; \textbf{interfaces conversacionales} para el propio pipeline (e.g.\ \texttt{telegram.bot} que responde con análisis al recibir un comando). \textbf{Consideraciones éticas y operativas:} costo por \textit{token}, privacidad de los datos enviados a APIs externas, no-determinismo de las respuestas (\texttt{temperature}, \textit{seed}), y la regla práctica: \orange{si una expresión regular resuelve el problema, no uses un LLM}.
\end{itemize}

\medskip
\noindent\textbf{Objetivos de aprendizaje}
\begin{enumerate}[leftmargin=*, itemsep=1pt]
    \item Justificar el patrón de integrar LLMs al pipeline de datos frente a usarlos como herramienta externa.
    \item Distinguir el uso de \texttt{httr2} (manual, transparente) frente al de \texttt{ellmer} (abstracción amigable) y elegir cuándo cada uno.
    \item Diseñar y consumir salidas estructuradas de un LLM con \textit{schemas} para integrar directamente al dataframe.
    \item Aplicar un LLM a tareas típicas de proyectos de datos: extracción estructurada, clasificación, resumen, normalización, etiquetado.
    \item Reconocer el panorama de extensiones (modelos locales, OCR, PII, interfaces conversacionales) y las consideraciones éticas y operativas de uso.
\end{enumerate}

\medskip
\noindent\textbf{Lecturas}
\begin{itemize}[leftmargin=*, itemsep=1pt]
    \item Documentación oficial de \texttt{ellmer} ---viñetas \textit{Getting started} y \textit{Structured data} \cite{wickham-ellmer}.
    \item Documentación oficial de \texttt{httr2} ---viñetas \textit{Getting started} y \textit{Wrapping APIs} \cite{wickham-httr2}.
    \item Documentación del proveedor LLM elegido (Anthropic, OpenAI, Google Gemini, Groq, etc.) ---referencia del \textit{endpoint} de \textit{chat completions} y de \textit{tool use} / \textit{structured output}.
\end{itemize}

% ==============================================================================
\section{Proyecto integrador}
% ==============================================================================

El curso se articula alrededor de un \blue{proyecto integrador} que cada estudiante construye de forma \textbf{individual} a lo largo del semestre, replicando el flujo completo de un proyecto de datos real.

\subsection*{Dataset y restricciones}

Cada estudiante elige una encuesta del \blue{INEGI} sujeta a tres restricciones:

\begin{enumerate}[leftmargin=*, itemsep=1pt]
    \item Debe contener \textbf{microdatos} (no agregados ni tabulados pre-procesados).
    \item Debe tener \textbf{al menos dos módulos relacionables} (e.g.\ ENSU con CS/CB/VIV; ENIGH con viviendas/hogares/personas/concentrado), para forzar \textit{joins} reales.
    \item Debe tener \textbf{al menos dos levantamientos en el tiempo} (panel, trimestral o varios años), para justificar el manejo de fechas y la acumulación.
\end{enumerate}

Encuestas que satisfacen estas restricciones incluyen ENSU, ENOE, ENIGH, ENVIPE, ENCIG, ENIF, ENADID y MCS-ENIGH, entre otras.

\subsection*{Pipeline}

El proyecto sigue el patrón \orange{ETL + análisis + producto}:

\begin{enumerate}[leftmargin=*, itemsep=1pt]
    \item \textbf{Extract.} Descargar los archivos de microdatos del sitio del INEGI; leer los descriptores técnicos.
    \item \textbf{Transform.} Corregir tipos, recodificar categóricas, manejar NAs, hacer \textit{joins} entre módulos, llevar a formato \textit{tidy}.
    \item \textbf{Load.} Guardar los datos transformados en un formato eficiente y reproducible (\texttt{.rds} y, hacia el final del semestre, base SQLite local).
    \item \textbf{Analizar.} Aplicar modelos estadísticos apropiados a las preguntas planteadas.
    \item \textbf{Producto.} \blue{Tablero Shiny} con visualizaciones interactivas más \blue{reporte periódico automatizado} con \texttt{officer}.
\end{enumerate}

\subsection*{Checkpoints}

El proyecto se evalúa por cinco \textit{checkpoints} alineados con las cinco partes del curso, más la entrega final:

\begin{center}
\begin{tabular}{cp{2.6cm}p{9cm}}
\toprule
\textbf{\#} & \textbf{Cierre de} & \textbf{Entregable} \\
\midrule
1 & Parte I (Sem.\ 4)   & Repositorio Git con estructura \texttt{files/}, \texttt{output/}, \texttt{docs/}, \texttt{pre/}; encuesta justificada (1 página); descriptor leído; primer script de extracción. \\
2 & Parte II (Sem.\ 7)  & ETL completo en \textit{scripts}: \textit{extract} + \textit{transform} + \textit{load} a \texttt{.rds}. Datos \textit{tidy} y validados. \\
3 & Parte III (Sem.\ 10) & Refactor del ETL en funciones reutilizables (\texttt{code/funciones.R}); \textit{master script} con \texttt{source()}. \\
4 & Parte IV (Sem.\ 12) & Análisis estadístico aplicado; \textit{load} conectado a base SQLite local. \\
5 & Parte V (Sem.\ 16) & Tablero Shiny funcional más reporte periódico automatizado con \texttt{officer}. \\
\bottomrule
\end{tabular}
\end{center}

% ==============================================================================
\section{Referencias principales}
% ==============================================================================

Las lecturas específicas por sesión se indican en el temario detallado. Las obras de consulta general son:

\nocite{wickham2023r4ds, wickham2019advr, wickham2016ggplot2, wickham2021mastering-shiny, rcoreteam-intro, bryan-happygit, ushey-renv, silge-tidytext, croissant-plm, robinson-broom, arelbundock-modelsummary, gohel-officer, gohel-flextable, bryan-googledrive, hester-gmailr, bslib-docs, kunst-highcharter, sievert-plotly, wickham-ellmer, wickham-httr2}

\printbibliography[heading=none]

\end{document}
